\documentclass[11pt,letterpaper]{article}


\usepackage[breaklinks, pagebackref=true, urlcolor=blue, colorlinks, citecolor=green!50!black, linkcolor=blue,bookmarksopen,bookmarksnumbered,linktocpage=true]{hyperref}

\usepackage[letterpaper, left=1in, right=1in, top=0.9in, bottom=0.9in]{geometry}
\usepackage[utf8]{inputenc}
\usepackage[american]{babel}
\usepackage[normalem]{ulem}
\usepackage{amsmath, amssymb, cases, amsthm}
\usepackage{thmtools}
\usepackage[shortlabels]{enumitem}
\usepackage[framemethod=TikZ]{mdframed}
\usepackage{bbm}
\usepackage{bm}
\usepackage{microtype}
\usepackage{xcolor}
\usepackage{makecell}
\usepackage{mathtools}
\usepackage{float}
\usepackage{modletters}
\usepackage{comment}
\usepackage{multirow}
\usepackage[sort]{cite}
%\usepackage{natbib}
\usepackage[capitalize,noabbrev]{cleveref}
\usepackage{graphics}
\usepackage{xspace}
% pdfsync omitted: editor synchronization only
\usepackage{tocloft}
\setlength\cftbeforesecskip{0.6em}

\usepackage{footnotehyper}
\makesavenoteenv{tabular}

\usepackage{pifont}
\newcommand{\cmark}{\ding{51}}%
\newcommand{\xmark}{\ding{55}}%

\providecommand{\mleftright}{}\let\mleft\left\let\mright\right

\declaretheorem[numberwithin=section,refname={Theorem,Theorems},Refname={Theorem,Theorems}]{theorem}
\declaretheorem[numberwithin=section,refname={Theorem,Theorems},Refname={Theorem,Theorems}]{thm}
\declaretheorem[numberlike=theorem]{lemma}
\declaretheorem[numberlike=theorem]{proposition}
\declaretheorem[numberlike=theorem]{corollary}
\declaretheorem[numberlike=theorem,style=definition]{definition}
\declaretheorem[numberlike=theorem, refname={Claim, Claims}, Refname={Claim, Claims}]{claim}
\declaretheorem[numberlike=theorem,style=remark]{remark}
\declaretheorem[numberlike=theorem,refname={Fact,Facts},Refname={Fact,Facts},name={Fact}]{fact}
\declaretheorem[numberlike=theorem, refname={Question,Questions},Refname={Question,Questions},name={Question}]{question}
\declaretheorem[numberlike=theorem, refname={Problem,Problems},Refname={Problem,Problems},name={Problem}]{problem}
\declaretheorem[numberlike=theorem, refname={Observation,Observations},Refname={Observation,Observations},name={Observation}]{observation}
\declaretheorem[numberlike=theorem, refname={Experiment,Experiments},Refname={Experiment,Experiments},name={Experiment}]{experiment}
\let\primitive\relax
\declaretheorem[numberlike=theorem,refname={Primitive,Primitives},Refname={Primitive,Primitives},name={Primitive}]{primitive}
\declaretheorem[numberlike=theorem,refname={Technique,Techniques},Refname={Technique,Techniques},name={Technique}]{technique}
\declaretheorem[numberlike=theorem,refname={Conjecture,Conjectures},Refname={Conjecture,Conjectures},name={Conjecture}]{conjecture}

\def\cameraready{0}  
\def\final{0}  
\def\iflong{\iffalse}
\ifnum\final=0  

\newcommand{\tawei}[1]{{\color{blue} [{\bf Ta-Wei:} #1]}}
\newcommand{\parth}[1]{{\color{purple} [{\bf Parth:} #1]}}
\newcommand{\TODO}[1]{{\color{blue!50!black} [{\bf Todo:} #1]}}
\else 
\newcommand{\yonggang}[1]{}
\newcommand{\danupon}[1]{}
\newcommand{\sagnik}[1]{}
\newcommand{\todo}[1]{}
\newcommand{\yuval}[1]{}
\newcommand{\jan}[1]{}
\newcommand{\blikstad}[1]{}
\newcommand{\tawei}[1]{}
\newcommand{\parth}[1]{}
\newcommand{\TODO}[1]{}
\fi  

\bibliographystyle{alpha}

\DeclarePairedDelimiter{\ceil}{\lceil}{\rceil}
\newcommand{\eps}{\varepsilon}

\newcommand{\mindeg}{\mathrm{mindeg}}
\newcommand{\nbh}{\mathrm{nbh}}
\newcommand{\spf}{\mathrm{spf}}
\newcommand{\direct}{MDCP }
\newcommand{\best}{\mathsf{best}}
\newcommand{\polylog}{\mathrm{polylog}}
\newcommand{\poly}{\mathrm{poly}}
\newcommand{\outdeg}{\mathrm{outdeg}}
\newcommand{\ones}{\mathrm{ones}}
\DeclareMathOperator*{\var}{Var}
\newcommand{\euler}{\mathrm{e}}
\newcommand{\set}[2][ ]{\{#2 \ifthenelse{\equal{#1}{ }}{ }{~|~#1}\}}
\newcommand{\fvedge}{\mathsf{FindViolatingEdge}}
\newcommand{\XORQ}{\mathsf{XOR}}
\newcommand{\ANDQ}{\mathsf{AND}}

\newcommand{\cc}{\mathrm{cc}}
\newcommand{\Rcal}{\mathcal{R}}
\newcommand{\Pcal}{\mathcal{P}}
\newcommand{\Acal}{\mathcal{A}}
\newcommand{\Scal}{\mathcal{S}}
\newcommand{\Qcal}{\mathcal{Q}}
\newcommand{\pflag}{\mathrm{pflag}}
\newcommand{\R}{\mathbb{R}}
\newcommand{\N}{\mathbb{N}}
\newcommand{\one}{\mathbf{1}}
\newcommand{\inact}{\mathsf{Inactive}}
\newcommand{\act}{\mathsf{Active}}
\newcommand{\cut}{\mathrm{cut}}
\newcommand{\cutd}{\overrightarrow{\mathrm{cut}}}
\newcommand{\Ed}{\overrightarrow{E}}
\newcommand{\mincut}{\lambda}
\newcommand{\tOh}{\widetilde{O}}
\newcommand{\ith}{i^{\scriptsize \mbox{{\rm th}}}}
\newcommand{\jth}{j^{\scriptsize \mbox{{\rm th}}}}



\newcommand{\intpot}[1]{Int_{pot}(#1)}


\newcommand{\ip}[1]{\mleft\langle #1\mright\rangle }
\newcommand{\expt}[2]{\underset{#1}{\mathbb{E}}\mleft[#2\mright]}
\newcommand{\A}{{\mathcal{A}}}
\newcommand{\BPM}{\mathsf{BPM}}
\newcommand{\BMM}{\mathsf{BMM}}
\newcommand{\UBPM}{\mathsf{UBPM}}

\newcommand{\T}{{\mathcal{T}}}

\newcommand{\U}{{\mathcal{U}}}

\newcommand{\bcc}[1]{\textbf{BCAST($#1$)}}
\newcommand{\dw}[1]{#1^{\downarrow}}

\newcommand{\MM}{\mathcal{P}}
\newcommand{\ORQ}{\mathsf{OR}}
\newcommand{\ISQ}{\mathsf{IS}}
\newcommand{\vol}{\mathrm{vol}}
\newcommand{\HPC}{\mathsf{HPC}}
\newcommand{\distance}{\mathsf{dist}}

\newcommand{\rank}{\mathsf{rank}}
\newcommand{\sspan}{\mathsf{span}}
\DeclareMathOperator*{\E}{\mathbb{E}}


\usepackage[most]{tcolorbox}
\tcbuselibrary{skins,breakable}
\tcbset{enhanced jigsaw}
\newcounter{myexample}
\usepackage{xparse}
\usepackage{lipsum}

\makeatletter
\newcommand\footnoteref[1]{\protected@xdef\@thefnmark{\ref{#1}}\@footnotemark}
\makeatother


\renewcommand{\paragraph}[1]{\medskip\noindent{\bf #1}\xspace}






\newcommand{\ground}{U}


\usepackage{tikz}
\usetikzlibrary{arrows}
\usetikzlibrary{arrows.meta}
\usetikzlibrary{shapes}
\usetikzlibrary{backgrounds}
\usetikzlibrary{positioning}
\usetikzlibrary{decorations.markings}
\usetikzlibrary{patterns}
\usetikzlibrary{calc}
\usetikzlibrary{fit}
\usetikzlibrary{decorations}


\renewcommand{\geq}{\geqslant}
\renewcommand{\leq}{\leqslant}
\renewcommand{\le}{\leq}
\renewcommand{\ge}{\geq}
% Parth's macros

\theoremstyle{definition}
\newtheorem{algocf}{Algorithm}
\newenvironment{Algorithm}{\begin{abox}\begin{algocf}}{\end{algocf}\end{abox}}

\Crefname{algocf}{Algorithm}{Algorithms}

\newenvironment{abox}{\begin{tcolorbox}[
		enlarge top by=5pt,
		enlarge bottom by=5pt,
		breakable,
		frame hidden,
		overlay broken = {
			\draw[line width=1pt, black]
			(frame.north west) rectangle (frame.south east);},
		overlay = {
			\draw[line width=1pt, black]
			(frame.north west) rectangle (frame.south east);},
		boxsep=0pt,
		left=4pt,
		right=4pt,
		top=10pt,
		arc=0pt,
		boxrule=1pt,toprule=1pt,
		colback=white
	]}
{\end{tcolorbox}}

\newcommand{\degen}{\textnormal{\textsf{DEGEN}}}

\newcommand{\fail}{\textsc{Reject}}
\newcommand{\true}{\textsc{Accept}}



\providecommand{\given}

\DeclarePairedDelimiter{\pparen}()
\DeclarePairedDelimiter{\card}\vert\vert

\DeclarePairedDelimiterX{\brac}[1][]{
%\renewcommand\given{\nonscript\: \delimsize\vert{} \nonscript\: \mathopen{}}
\ifblank{#1}{\,\cdot\,}{#1}}

\DeclarePairedDelimiterX{\Set}[1]\{\}{
  \renewcommand\given{\nonscript\: \delimsize\vert{} \nonscript\: \mathopen{}}
\ifblank{#1}{\,\cdot\,}{#1}}

\renewcommand{\given}{\nonscript\: \vert \nonscript\:}


\NewDocumentCommand{\Prob}{sO{}E{_}{{}}m}{%
    \Pr_{#3}
    \IfBooleanTF{#1}
    {\brac*{#4}}
    {\brac[#2]{#4}}
}

\DeclarePairedDelimiterXPP{\Ot}[1]{\widetilde{O}}(){}{#1}
\DeclarePairedDelimiterXPP{\Omgt}[1]{\widetilde{\Omega}}(){}{#1}

\DeclarePairedDelimiterXPP{\Var}[1]{\var}[]{}{#1}

\newcommand{\dsi}{\cD_\SI}
\newcommand{\pisi}{\pi_\SI}

\newcommand{\ptd}{\Lambda}
\DeclarePairedDelimiterXPP{\PTD}[2]{\Lambda}(){}{#1, #2}

\DeclareMathOperator*{\distrib}{dist}
\NewDocumentCommand{\Dist}{sO{}E{_}{{}}m}{%
  \mu_{#3}
  \IfBooleanTF{#1}
  {\pparen*{#4}}
  {\pparen[#2]{#4}}
}

\DeclareMathOperator*{\expect}{\mathbb{E}}
\NewDocumentCommand{\Exp}{sO{}E{_}{{}}m}{%
  \expect_{#3}
  \IfBooleanTF{#1}
  {\brac*{#4}}
  {\brac[#2]{#4}}
}

\newcommand{\re}{\rv{e}}

\newcommand{\estar}{e^{\ast}}
\newcommand{\restar}{\re^{\ast}}

\newcommand{\rPi}{\rv{\Pi}}

\newcommand{\Perm}{\Gamma}
\newcommand{\rPerm}{\rv{\Gamma}}

\newcommand{\uni}{\mathcal{U}}

\newcommand{\qpi}{q^{\Pi}}

\newcommand{\score}{\mathsf{Score}}
\newcommand{\totscore}{\mathsf{TScore}}


\DeclareMathOperator*{\inform}{\mathbb{I}}
\DeclarePairedDelimiterXPP{\Inf}[1]{\inform}(){}{#1}

\DeclareMathOperator*{\entropy}{\mathbb{H}}
\DeclarePairedDelimiterXPP{\Ent}[1]{\entropy}(){}{#1}

\DeclareMathOperator{\myIC}{\mathsf{IC}}
\DeclareMathOperator{\myCC}{\mathsf{CC}}
\newcommand{\ExtIC}{\myIC^{\textnormal{ext}}}
\newcommand{\IntIC}{\myIC^{\textnormal{int}}}

\newcommand{\rPTD}{\mathbf{\Lambda}}

\newcommand{\Leq}[1]{{\underset{\textnormal{#1}}\leq}}
\newcommand{\Geq}[1]{{\underset{\textnormal{#1}}\geq}}
\newcommand{\Eq}[1]{{\underset{\textnormal{#1}}=}}



\newcommand{\bhpc}{\ensuremath{\textnormal{\textsf{BHPC}}}\xspace}
\newcommand{\bmhpc}{\ensuremath{\textnormal{\textsf{BMHPC}}}\xspace}
\newcommand{\hpc}{\ensuremath{\textnormal{\textsf{HPC}}}\xspace}
\newcommand{\mhpc}{\ensuremath{\textnormal{\textsf{MHPC}}}\xspace}
\newcommand{\setint}{\ensuremath{\textnormal{\textsf{SetInt}}}\xspace}
\newcommand{\defc}{\mathtt{def}}
\newcommand{\odeg}{\text{odeg}}
\newcommand{\Vaux}{V_{\text{aux}}}
\newcommand{\Eaux}{E_{\text{aux}}}
\newcommand{\prantar}[1]{{\color{blue} [{\bf Prantar:} #1]}}
\newcommand{\mypar}[1]{\smallskip\noindent{\sffamily\bfseries #1.}~}
\newcommand{\floor}[1]{\lfloor #1 \rfloor}




\newcommand{\cmax}{\ensuremath{\textnormal{\textsf{c}}_{\textnormal{rs}}}}
\newcommand{\protMO}{\ensuremath{\Prot_{\mo}}\xspace}
\newcommand{\bProtMO}{\ensuremath{\bm{\Prot}_{\mo}}\xspace}
\newcommand{\protBHH}{\ensuremath{\Prot_{\bhh}}\xspace}
\newcommand{\bProtBHH}{\ensuremath{\bm{\Prot}_{\bhh}}\xspace}


\newcommand{\Paren}[1]{\Big(#1\Big)}
\newcommand{\Bracket}[1]{\Big[#1\Big]}
\newcommand{\bracket}[1]{\mleft[#1\mright]}
\newcommand{\paren}[1]{\ensuremath{\mleft(#1\mright)}\xspace}

\newenvironment{tbox}{\begin{tcolorbox}[
		enlarge top by=5pt,
		enlarge bottom by=5pt,
		 breakable,
		 boxsep=0pt,
                  left=4pt,
                  right=4pt,
                  top=10pt,
                  boxrule=1pt,toprule=1pt,
                  colback=white,
                  arc=-1pt,
        %          drop shadow={black,opacity=1},
                 % arc=0pt
                  ]%%
	}
{\end{tcolorbox}}



\newcommand{\Prot}{\ensuremath{\Pi}}
\newcommand{\prot}{\ensuremath{\pi}}
\newcommand{\Rpub}{{R}^{pub}}
\newcommand{\Rpri}{{R}^{pri}}

\newcommand{\bProt}{\bm{\Prot}}
\newcommand{\bRpub}{\bm{\Rpub}}
\newcommand{\bRpri}{\bm{\Rpri}}


\newcommand{\bV}{\bm{V}}
\newcommand{\bX}{\bm{X}}
\newcommand{\bY}{\bm{Y}}
\newcommand{\bZ}{\bm{Z}}
\newcommand{\bA}{\bm{A}}
\newcommand{\bP}{\bm{P}}
\newcommand{\bR}{\bm{R}}
\newcommand{\bM}{\bm{M}}
\newcommand{\bI}{\bm{I}}
\newcommand{\bJ}{\bm{J}}
\newcommand{\bSigma}{\bm{\sigma}}
\newcommand{\bG}{\bm{G}}

\newcommand{\bB}{\ensuremath{\bm{B}}}
\newcommand{\bC}{\ensuremath{\bm{C}}}
\newcommand{\bD}{\ensuremath{\bm{D}}}

\newcommand{\supp}[1]{\ensuremath{\textsc{supp}(#1)}}




\newtheorem{mdresult}{Result}
\newenvironment{result}{\begin{mdframed}[backgroundcolor=lightgray!40,topline=false,rightline=false,leftline=false,bottomline=false,innertopmargin=2pt]\begin{mdresult}}{\end{mdresult}\end{mdframed}}


%-------------------------- Textbox ---------------------------------

\newcommand{\textbox}[2]{
{
\begin{tbox}
\textbf{#1}
{#2}
\end{tbox}
}
}

\newcommand{\rA}{\rv{A}}
\newcommand{\rB}{\rv{B}}
\newcommand{\rC}{\rv{C}}
\newcommand{\rD}{\rv{D}}

\newcommand{\rX}{\rv{X}}
\newcommand{\rY}{\rv{Y}}
\newcommand{\rI}{\rv{I}}

\newcommand{\rv}[1]{\ensuremath{\mathsf{#1}}}

\newcommand{\II}{\ensuremath{\mathbb{I}}}
\newcommand{\HH}{\ensuremath{\mathbb{H}}}
\newcommand{\mi}[2]{\ensuremath{\II(#1 \,; #2)}}
\newcommand{\en}[1]{\ensuremath{\HH(#1)}}

\newcommand{\itfacts}[1]{Fact~\ref{fact:it-facts}-(\ref{part:#1})\xspace}

\newcommand{\DD}{\ensuremath{\mathbb{D}}}

\newcommand{\distribution}[1]{\ensuremath{\textnormal{dist}(#1)}\xspace}

\newcommand{\XX}{\ensuremath{\mathcal{X}}}
\newcommand{\YY}{\ensuremath{\mathcal{Y}}}


\newcommand{\HPCk}{\ensuremath{\HPC_k}\xspace}

\newcommand{\protHPC}{\ensuremath{\prot_{\HPC}}}

\newcommand{\distHPC}{\ensuremath{\dist_{\HPC}}}

\newcommand{\unif}{\ensuremath{\mathcal{U}}}

\newcommand{\rZ}{\rv{Z}}
\newcommand{\rt}{\rv{t}}
\newcommand{\rT}{\rv{T}}
\newcommand{\rE}{\rv{E}}
\newcommand{\rP}{\rv{P}}
\newcommand{\rProt}{\rv{\Prot}}

\newcommand{\distSI}{\ensuremath{\dist}_{\ensuremath{\textnormal{\textsf{SI}}}}\xspace}


\newcommand{\protSI}{\ensuremath{\prot}_{\ensuremath{\textnormal{\textsf{SI}}}}\xspace}
\newcommand{\ProtSI}{\ensuremath{\Prot}_{\ensuremath{\textnormal{\textsf{SI}}}}\xspace}
\newcommand{\rProtSI}{\ensuremath{\rv{\Prot}}_{\ensuremath{\textnormal{\textsf{SI}}}}\xspace}

\newcommand{\distPI}{\ensuremath{\dist}_{\ensuremath{\textnormal{\textsf{PI}}}}\xspace}
\newcommand{\protPI}{\ensuremath{\prot}_{\ensuremath{\textnormal{\textsf{PI}}}}\xspace}
\newcommand{\ProtPI}{\ensuremath{\Prot}_{\ensuremath{\textnormal{\textsf{PI}}}}\xspace}
\newcommand{\rProtPI}{\ensuremath{\rv{\Prot}}_{\ensuremath{\textnormal{\textsf{PI}}}}\xspace}


\newcommand{\SI}{\ensuremath{\textnormal{\textsf{Set-Int}}}\xspace}
\newcommand{\PI}{\ensuremath{\textnormal{\textsf{Pair-Int}}}\xspace}


\newcommand{\errs}{\ensuremath{\textnormal{errs}}\xspace}




\newcommand{\rK}{\rv{K}}
\newcommand{\rJ}{\rv{J}}
\newcommand{\rL}{\rv{L}}
\newcommand{\rS}{\rv{S}}
\newcommand{\rM}{\rv{M}}
\newcommand{\rR}{\rv{R}}
\newcommand{\rbarS}{\rv{\barS}}



\newcommand{\distDisjN}{\ensuremath{\distDisj^{\textnormal{\textsf{N}}}}\xspace}
\newcommand{\distDisjY}{\ensuremath{\distDisj^{\textnormal{\textsf{Y}}}}\xspace}

\newcommand{\mii}[3]{\ensuremath{\II_{#3}(#1 \,; #2)}}

\newcommand{\disjin}[4]{ \ensuremath{\begin{bmatrix} #1&\hspace{-8pt}#2 \\ #3&\hspace{-10pt}#4 \end{bmatrix}}}
\renewcommand{\disjin}[4]{\textnormal{[\ensuremath{#1#2,#3#4}]}}

\newcommand{\barS}{\overline{S}}

\newcommand{\hd}[2]{\ensuremath{\textnormal{h}(#1,#2)}\xspace}
\newcommand{\hdt}[2]{\ensuremath{\textnormal{h}^2(#1,#2)}\xspace}


\newcommand{\CC}[2]{\ensuremath{\textnormal{\textsf{CC}}_{#2}(#1)}\xspace}
\newcommand{\ICost}[2]{\ensuremath{\textnormal{\textsf{IC}}_{#2}(#1)}\xspace}
\newcommand{\IC}[2]{\ICost{#1}{#2}}


\newcommand{\rTheta}{\rv{\Theta}}



\newcommand{\comp}{\succ_{\ProtSI}}
\newcommand{\compr}{\prec_{\ProtSI}}
\newcommand{\factor}{3}
%%%%
\title{Optimal round-communication lower bound for Hidden Pointer Chasing}
\author{Sepehr Assadi, Prantar Ghosh, Bruno Loff, Parth Mittal, Sagnik Mukhopadhyay}
\date{doi:10.4230/LIPIcs.CCC.2024.7; CC BY 4.0}
\begin{document}
\maketitle
\setcounter{section}{2}
\noindent\textit{Editorial note.} Native fixed arXiv v1 author TeX under CC BY 4.0. Counted result: thm:lb-hpc. The complete common lower-bound section is retained so every supporting epsilon-solver argument and cross-reference remains available; the other stated main result is supporting context and is not counted in this record. Appendix input directives are mechanically expanded via original-file spans. Unrelated streaming reduction appendix is omitted.\par\medskip
\noindent\textit{Allowed standard prerequisites (editorial).}\begin{itemize}
\item Set-intersection internal information lower bound. Jayram, Kumar and Sivakumar 2003; author Proposition prop:SI-hard and exact distribution retained
\item Standard entropy, mutual information, rectangle property, Chernoff and Pinsker facts. Author information-theory and sampling appendix retained; Cover and Thomas 2006 and cited sources
\end{itemize}\medskip
\subsection{\texorpdfstring{$k$-Cores}{k-Cores} and Degeneracy in Graph Streams}\label{sec:degen-intro}

For any undirected graph $G=(V,E)$ and integer $k \geq 1$, a \emph{$k$-core} in $G$ is a maximal set $S$ of vertices such that the induced subgraph of $G$ on $S$, denoted by $G[S]$, has a minimum degree of at least $k$. In other words, 
any vertex in $S$ has at least $k$ other neighbors in $S$. 

$k$-Cores provide a natural notion of well-connectedness in massive graphs, and as such, computing $k$-cores (and more generally $k$-core decompositions; see, e.g.,~\cite{LiuSYDS22}) 
has been widely studied in databases~\cite{BonchiGKV14, ChuZ00ZXZ20, LiZZQZL19}, social networks~\cite{DhulipalaBS17, DhulipalaBS18, KhaouidBST15}, machine learning~\cite{Alvarez-HamelinDBV05, EsfandiariLM18, GhaffariLM19}, 
among others~\cite{LiWSX18,GalimbertiBGL20, SariyuceGJWC13}. 

As a result, in recent years, there has been a rapidly growing body of work on
computing $k$-cores on massive graphs in parallel and streaming
models of computation~\cite{SariyuceGJWC13,DhulipalaBS17, DhulipalaBS18,EsfandiariLM18, GhaffariLM19,LiuSYDS22}. In particular,~\cite{EsfandiariLM18} presented a single-pass algorithm that for any $\eps > 0$, 
computes a $(1-\eps)$-approximation of every $k$-core in $G$ (i.e., obtains a $(1-\eps)$-approximate $k$-core decomposition) in $\Ot{n/\eps^2}$ space (see also~\cite{SariyuceGJWC13} for an earlier streaming algorithm and~\cite{GhaffariLM19}
for a closely related parallel algorithm). 

The \emph{degeneracy} of a graph $G=(V,E)$, denoted by $\kappa(G)$, is the largest integer $k \geq 0$ such that $G$ contains a non-empty $k$-core. 
The simple greedy algorithm that at every step peels off the smallest degree vertex results in the so-called \emph{degeneracy ordering} of $G$ and $\kappa{(G)}$ 
is equal to the largest degree of a vertex removed in this peeling process~\cite{MatulaB83}. Degeneracy is a standard measure of uniform sparsity and is closely
related to other such notions like arboricity (which is within a factor $2$ of degeneracy). Moreover, computing degeneracy is a subroutine for approximating various
other problems such as arboricity~\cite{AlonYZ97}, densest subgraph~\cite{Charikar00}, $(\kappa+1)$ coloring~\cite{erdHos1966chromatic}. 

The degeneracy problem, and closely related uniform-sparsity measures such as densest subgraph, have also been studied extensively in the graph 
streaming literature~\cite{BahmaniKV12,Farach-ColtonT14,BhattacharyaHNT15,McGregorTVV15,Farach-ColtonT16,BeraCG19,AlonA20}. In particular,~\cite{Farach-ColtonT14} provided an $O(\log{n})$-pass semi-streaming algorithm that outputs a constant factor approximation to degeneracy, and~\cite{Farach-ColtonT16} subsequently improved this to a single-pass $(1-\eps)$-approximation in $\Ot{n/\eps^2}$ space (see also~\cite{McGregorTVV15} for densest subgraph and~\cite{BeraCG19,AlonA20} for degeneracy coloring). 
 
\medskip

In terms of lower bounds, \cite{BeraCG19} prove that any \emph{single-pass} streaming algorithm that computes the exact value of degeneracy or approximates it within an additive factor of $\lambda$ requires $\Omega(n^2)$ space or $\Omega(n^2/\lambda^2)$ space respectively. Our \emph{polynomial-pass} lower bounds for $k$-cores and degeneracy, now in a very strong sense, rule out the possibility of extending any prior semi-streaming algorithms computing near-optimal solutions to these problems, to compute \emph{exactly} optimal solutions. 


\section{Preliminaries} \label{sec:prelim}



\subsection{The Set-Intersection Problem}

\begin{definition}
    The \textbf{$m$-bit Set-Intersection} problem is a communication problem
    for two players Alice and Bob.
    Alice is given a subset $X \subseteq [m]$ and Bob is given a subset
    $Y \subseteq [m]$, with the promise that $|X \cap Y| = 1$.
    Their goal is to learn the point of intersection $\{t\} = X \cap Y$. 
\end{definition}

Let us define a distribution $\cD_\SI$ over inputs for Set-Intersection. This is the positive case of Razborov's original hard distribution for disjointness \cite{Razborov92}.

\begin{tbox}
    \textbf{Distribution $\cD_\SI$}~\cite{AssadiR20} on sets
    $(X, Y)$ from the universe $\cX=[m]$: 

    \begin{itemize}
    \item Uniformly sample two disjoint sets $X', Y' \subseteq [m]$ of size
      $\frac m 4 - 1$.
    \item Sample $\estar \in [m]\setminus(X' \cup Y')$ uniformly at random
      and define $X = X' \cup \Set{\estar}$ and $Y = Y' \cup \Set{\estar}$.
\end{itemize}
\end{tbox}

\begin{proposition}\label{prop:SI-hard}
  The internal information complexity of \emph{finding} the intersecting element
  in a $\setint_m$ instance sampled from $\dsi$ is $\Omega(m)$.
\end{proposition}

\Cref{prop:SI-hard} is not directly implied by standard set-intersection lower
bounds, since they use distributions where $X$ and $Y$ never intersect.
To obtain this proposition, one can (for example) apply the result
of~\cite{JayramKS03} on the information cost of set disjointness protocols on
intersecting distributions (see~\cite{AssadiCK19} for more details).

\subsection{The (Boolean) Hidden Pointer Chasing Problem}

The Hidden Pointer Chasing ($\hpc$) problem is a composition of Pointer Chasing and Set-Intersection, defined as follows.

\begin{definition}\label{prob:hpc}
The \textbf{Hidden Pointer Chasing ($\HPC$)} problem is a communication problem
with four players $P_A, P_B, P_C$, and $P_D$. 
Let $\cX := \{x_1,...,x_m\}$ and $\cY := \{y_1,...,y_m\}$ be two disjoint universes. 
    \begin{itemize}
        \item For any $x\in \cX$, $P_A$ and $P_B$ are given an instance $(A_x,B_x)$ of \setint over the universe $\cY$ where $A_x \cap B_x = \{t_x\}$ for a single target element $t_x \in \cY$. We define $A := \{A_{x_1} ,\ldots, A_{x_m} \}$ and $B:=\{B_{x_1} ,\ldots, B_{x_m} \}$ as the whole input to $P_A$ and $P_B$, respectively.
        
        \item For any $y\in \cY$, $P_C$ and $P_D$ are given an instance $(C_y,D_y)$ of $\setint$ over the universe $\cX$ where $C_y \cap D_y = \{t_y\}$ for a single target element $t_y \in \cX$. We define $C := \{C_{y_1} ,\ldots, C_{y_m} \}$ and $D:=\{D_{y_1} ,\ldots,D_{y_m} \}$ as the whole input to $P_C$ and $P_D$, respectively.
        
        \item We define two mappings $f_{AB} : \cX \to \cY$ and $f_{CD} : \cY \to \cX$ such that:
        \[
            (a) ~ \textnormal{for any $x\in \cX$, $f_{AB}(x) = t_x$;} \qquad \qquad (b) ~ \textnormal{for any $y\in \cY$, $f_{CD}(y) = t_y$}.
         \]
        \item Let $x_1 \in \cX$ be an arbitrary fixed element known to all players. The pointers $z_0,z_1,z_2,\ldots$ are defined inductively as 
        $$z_0 := x_1, ~~z_1 := f_{AB}(z_0),~~ z_2 := f_{CD}(z_1), ~~z_3 := f_{AB}(z_2),~\ldots.$$
    \end{itemize}
For any integers $m,r \geq 1$, the \textnormal{\textbf{$m$-bit, $r$-step Hidden Pointer Chasing}} problem, denoted by $\hpc_{m,r}$ is defined as the communication problem of finding the pointer $z_r$. Meaning, we wish to output the index $i \in [m]$ such that $z_r = x_i$ (when $r$ is even) or $z_r = y_i$ (when $r$ is odd). 

We also define a Boolean variant of $\hpc$ as follows. In the \textnormal{\textbf{$m$-bit, $r$-step \underline{Boolean} Hidden Pointer Chasing}} problem, denoted by $\bhpc_{m,r}$, the goal is to output the bit $b(z_r):=i \mod 2$, where $z_r = x_i$ or $y_i$ depending on if $r$ is even or odd respectively. 
\end{definition}

\subsection{The (Boolean) Multilayer Hidden Pointer Chasing Problem}

The Multilayer Hidden Pointer Chasing problem ($\mhpc$) combines
$r$ instances of $\hpc$ as its input in the following manner.


\begin{definition}
  The \textbf{Multilayer Hidden Pointer Chasing} ($\mhpc$) problem is a
  communication problem with four players $P_A, P_B, P_C$ and $P_D$. Let $m,r$ be integers $\geq 1$. The $m$-bit, $r$-layer \mhpc problem, denoted by $\mhpc_{m,r}$ is defined as follows.
  \begin{itemize}
    \item For each $i \in \bracket{r}$, the players receive an $\hpc$ instance
      $(A^i, B^i, C^i, D^i)$, each over the disjoint universes
      $\cX$ and $\cY$ of size $m$.
      Let $t^i_x$ denote the unique element in the intersection
      $A^i_x \cap B^i_x$, and $t^i_y$ the unique element in $C^i_y \cap D^i_y$.
    \item For each $i \in \brac{r}$, we have the maps $f^i_{AB} : \cX \to \cY$
      and $f^i_{CD} : \cY \to \cX$, defined as:
      \[
        (a) ~ \textnormal{for any $x\in \cX$, $f^i_{AB}(x) = t^i_x$;}
        \qquad \qquad
        (b) ~ \textnormal{for any $y\in \cY$, $f^i_{CD}(y) = t^i_y$}.
      \]
    \item
      The pointers $z_i$ are defined inductively as:
      \[
        z_0 := x_1, ~~z_1 := f^1_{AB}(z_0),~~ z_2 := f^2_{CD}(z_1),
        ~~z_3 := f^3_{AB}(z_2),~\ldots.
      \]
  \end{itemize}
  % As before, the goal for the players is to compute $z_r$.

  The goal of the problem is to find $z_r$, i.e., the index $i \in [m]$ such that $z_r = x_i$ (when $r$ is even) or $z_r = y_i$ (when $r$ is odd). 

Analogous to \hpc, we define the Boolean variant of $\mhpc$ as follows. In the \textnormal{\textbf{$m$-bit, $r$-layer \underline{Boolean} \mhpc}} problem, denoted by $\bmhpc_{m,r}$, the goal is to output the bit $b(z_r):=i \mod 2$, where $z_r = x_i$ or $y_i$ depending on if $r$ is even or odd respectively. 
\end{definition}

We remark that $C^j, D^j$ for odd $j$, and $A^j, B^j$ for even $j$, are redundant (they do not affect the output of the function). They are kept for notational convenience and ease of presentation. 

\subsection{Misaligned \texorpdfstring{$r$-Round}{r-Round} Protocols}

We are interested in bounded-round protocols for $\hpc$, $\bhpc$, $\mhpc$ and $\bmhpc$. In each round, two of the players speak with each other ($P_A$ with $P_B$, or $P_C$ with $P_D$), but not with the other pair of players, with a message between pairs of players triggering the next round.

\begin{definition}[Round of $\hpc_{m,r}$]\label{def:round}
An $r$-round protocol is a protocol where the communication is organized in rounds $1, 2, \ldots, r$. A given pair of players --- $P_A$ and $P_B$ or $P_C$ and $P_D$ --- communicate with each other in the odd rounds (the first, third, etc rounds) and the other pair of players communicate with each other in the even rounds. Each round ends with the speaking pair of players sending a message to the other pair of players, and this marks the beginning of the next round.
\end{definition}

It is easy to see that one can solve any of the above HPC variants by $r$-round protocols communicating $O(r m)$ bits, if $P_A$ and $P_B$ speak in the odd rounds, and $P_C$ and $P_D$ speak in the even rounds. We will show lower-bounds for \textit{misaligned} $r$-round protocols where the \textit{wrong} pair of players speak:

\begin{definition}
We call an $r$-round protocol \textit{misaligned} if the pair of players $P_C$ and $P_D$ speak in the odd rounds, and $P_A$ and $P_B$ speak in the even rounds.
\end{definition}

The above definitions naturally suggest that, in an $r$-round protocol, we should think of $z_j$ as the pointer that the players wish to know by the end of round $j$. In a misaligned protocol, however, the players speaking in round $j$ do not hold the Set-Intersection instance that needs to be solved to find $z_j$.

\subsection{Notation and Terminology}

\mypar{Notation for Protocols} We use the notation $\prot_\SI$, $\prot_\bhpc$,
and $\prot_\bmhpc$, to denote protocols for Set-Intersection, $\bhpc$ and
$\bmhpc$, respectively. 

In general, we will use the upright $\rA$ to denote a random variable
associated with $A$.

Now consider a situation where we choose inputs $(A, B, C, D)$ from some
distribution on inputs to $\hpc_{m,r}$.
We then let $\rZ_1, \ldots , \rZ_r$ denote random variables with values
$z_1, \ldots, z_r$ of \Cref{prob:hpc}.
We use the variable $j \in [r]$ to index the rounds of a
$r$-round protocol.
For any $j \in [r]$, we define $\Prot_j$ as the set of all messages
communicated in round $j$, and more generally $\Prot$ denotes the entire
protocol transcript. 

For any $x \in \cX$ and $y \in \cY$, we define random variables
$\rT_x \in \cY$ and $\rT_y \in \cX$ giving the intersecting elements of the
Set-Intersection instance $(\rA_x, \rB_x)$ and $(\rC_y, \rD_y)$, respectively.
We define $\rE_j = (\rPi_1, \ldots, \rPi_{j-1}, \rZ_1, \ldots, \rZ_{j-1})$ to
include all messages exchanged before round $j$, together with the entire path
followed by the pointers before round $j$.
Note that during round $j$, the players wish to compute the intersection
$A_{z_j} \cap B_{z_j}$ (or $C_{z_j} \cap D_{z_j}$), so knowing $\rE_j$ should
not help them to do that. (Note that $\rE_1$ is empty.)

On an input $(A, B, C, D)$ to $\bmhpc_{m,r}$, we define the notation
$Z_i$, $j$, $\Prot_j$, $T^j_x$, $T^j_y$ analogously.

If we have random variables $\rA, \rB$, we will broadly use the notation
$\mu(\rA \mid \rB)$ to mean ``the distribution of $\rA$, conditioned on the
value of $\rB$ that occurred''.
So, for example, if $\rB$ is constant we may write $\mu(\rA)$, and this
is some fixed distribution.
But if $\rB$ is non-constant, the distribution $\mu(\rA \mid \rB)$ may vary
depending on $\rB$, and so $\mu(\rA \mid \rB)$ is itself also a random
variable, which depends on $\rB$ and can be determined by knowing $\rB$ alone
(even without knowing $\rA$).

\mypar{Other notation}
An input graph to a streaming algorithm will usually be called $G=(V,E)$ and $n$ will denote $|V|$. The degeneracy of graph $G$ will be denoted by $\kappa(G)$, or simply $\kappa$ if the graph is clear from the context. The notation ``$\log x$'' stands for $\log_2 x$. For a positive integer $a$, the notation $[a]$ denotes the set $\{1,2,\ldots,a\}$ and for integers $a<b$, $[a,b]$ denotes the set $\{a,a+1,\ldots,b\}$.
 


\subsection{Positive Triangular Discrimination} \label{sec:td}

\begin{definition}[Positive triangular discrimination]
    Given two distributions $\mu$ and $\nu$ on a sample space $\Omega$, the \textit{positive triangular discrimination} 
    % $\Delta(\mu, \nu)$ is defined as the following measure of distance between $\mu$ and $\nu$:
    % \[
    % \Delta(\mu, \nu) = \sum_{x \in x}\frac{(\mu(x) - \nu(x))^2}{\mu(x) + \nu(x)}.
    % \]
    % The positive triangular discrimination 
    $\Lambda(\mu,\nu)$ is:
    \[
    \Lambda(\mu, \nu) = \sum_{\substack{x \in \Omega\\\mu(x) > \nu(x)}}\frac{(\mu(x) - \nu(x))^2}{\mu(x) + \nu(x)}.  %\leq \Delta(\mu, \nu).
    \]
\end{definition}

In this paper, we work exclusively with \textit{positive} triangular discrimination. Below we state some of the useful properties of positive triangular discrimination. We defer the proofs of the following propositions to \Cref{sec:td-props}.


\begin{restatable}{proposition}{lambdatvd}\label{prop:prop-lambda}
$0 \le \frac{\|\mu - \nu\|^2_1} 8\leq \Lambda(\mu, \nu)\leq \frac{\|\mu - \nu\|_1} 2 \le 1$. 
\end{restatable}


The following property is already implicit in \cite{Yehudayoff20}, but making it explicit now simplifies the presentation later. Note that an analogous equation for total variation distance, i.e., $\E_\mu[f(x)] \le \|\mu-\nu\|_1\cdot \max_{x \in x} f(x) + \cdot \E_{\nu}[f(x)]$ can be obtained almost from the definition. The following statement, however, needs a bit more work.

\begin{restatable}{proposition}{lambdaloss}\label{itm:lambda-loss}
For any function $f:\Omega\to\bbR_{\ge 0}$,
\[ \E_\mu[f(x)] \le \Lambda(\mu,\nu)\cdot \max_{x \in \Omega} f(x) + 6 \cdot \E_{\nu}[f(x)]. \]
\end{restatable}



\noindent
We will also make use of the convexity of $\Lambda$.

\begin{restatable}{proposition}{convex}\label{prop:convex}
Positive triangular discrimination $\Lambda(\mu,\nu)$ is convex as a joint function of the two distributions $\mu$ and $\nu$. 
\end{restatable}



\subsection{Degeneracy and \texorpdfstring{$k$-Cores}{k-Cores}}

The definition of a $k$-core and the degeneracy $\kappa(G)$ of a graph $G$ can be found in \Cref{sec:degen-intro}. Here, we define some related terms that come up often in the paper. We also state some basic facts about them that we later use in our proofs. 


Given an ordering $\sigma$ of vertices, we say $v\prec_\sigma u$ if $u$ appears (somewhere) after $v$ in $\sigma$. 

\begin{definition}[Outdegree with respect to an ordering]\label{def:odegordering}
    For a graph $G=(V,E)$. given an ordering $\sigma$ of vertices in $V$, the outdegree of a vertex $v$ with respect to $\sigma$, denoted by $\odeg_{\sigma}(v)$, is the number of neighbors of $v$ that appear after it in $\sigma$. Formally,
$$\odeg_{\sigma}(v) = |\{u\in N(v): v \prec_\sigma u\}|$$ 
\end{definition}

We can also think of it as orienting all edges $(x,y)$ from $x$ to $y$ if $x\prec_\sigma y$. Then $\odeg_\sigma(v)$ is the outdegree of $v$ in the resultant oriented graph. 


\begin{definition}[$k$-ordering]\label{def:k-ordering} 
    Given a graph $G=(V,E)$, an ordering $\sigma$ of vertices in $V$ is called a $k$-ordering if $\forall v\in V: \odeg_\sigma(v)\leq k$.   
\end{definition}

The following fact connects vertex orderings to degeneracy. 

\begin{fact}[\cite{MatulaB83}]\label{fact:min-k-ord}
    The minimum number $k$ for which a graph $G$ admits a $k$-ordering is $\kappa(G)$.
\end{fact}


The above fact implies that to prove that a graph has degeneracy at most $k$, it suffices to exhibit a $k$-ordering of the nodes. We now formally define \emph{degeneracy ordering} of a graph.  

\begin{definition}[Degeneracy ordering] \label{def:degen-ord}
    Given a graph $G=(V,E)$, a $\kappa(G)$-ordering of $V$ is called a degeneracy ordering of $G$.
\end{definition}

A degeneracy ordering of a graph can be obtained by running the following ``peeling algorithm'' on it: starting with an empty sequence $\sigma$, recursively remove the minimum degree vertex from $G$ and append it to $\sigma$ until the graph becomes empty. Again, the maximum degree that a node has when it is peeled by this algorithm is exactly $\kappa(G)$. If not, then the peeling algorithm gives a $k$-ordering for some $k<\kappa(G)$ which violates \Cref{fact:min-k-ord}. Hence, we get the following fact. 


 
  

\begin{fact}\label{fact:peelingmin-degen}
Suppose the peeling algorithm on a graph $G$ removes its nodes in the order $v_1,\ldots,v_n$. Define $G_0:=G$ and $G_i:= G\setminus \{v_1,\ldots,v_i\}$ for $i\in [n]$. Then, $$\kappa(G)=\max_{i\in [n]} \deg_{G_{i-1}}(v_i)$$. 
\end{fact}





\newcommand{\bucket}[1]{\textsf{Bucket} #1}

\section{New Lower Bounds for HPC and Multilayer HPC} \label{sec:lower-bound-hpc}

The two main results of this section are an $\Omega(m^2)$ lower-bound against misaligned $r$-round protocols for $\bmhpc_{m,r}$, and an $\Omega(\frac{m^2}{r})$ lower-bound against misaligned $r$-round protocols for $\bhpc_{m,r}$. Notice that in a misaligned $r$-round round protocol for $\bmhpc_{m,r}$ or $\bhpc_{m,r}$, it is players $P_C$ and $P_D$ who begin the protocol by talking with each other. This means that the ``wrong'' pair of players begin to speak, in the sense that they wish to compute the value $\{z_1\} = A_{z_0} \cap B_{z_0}$, but this instance is with Alice and Bob, so they have no way to do this. So the first round of communication says nothing about $z_1$: the best $P_C$ and $P_D$ can do is send some information about all of their Set-Intersection instances, without knowing which is important. So each bit that $P_C$ and $P_D$ communicate with $P_A$ and $P_B$ in the first round, can only reveal $\frac 1 m$ bits of information about the average instance. So, in the next round, $P_A$ and $P_B$ cannot have learned what $C_{z_1}$ or $D_{z_1}$ are, either. But then, how can they learn $z_2$? The difficult situation is now reversed! This ``always one step behind'' situation is similar to what happens for pointer chasing \cite{NisanW93, PonzioRV99, Yehudayoff20}, except now the pointers are ``hidden'' behind set intersection instances.

\medskip\noindent
How much communication must the players use in order to escape this unfortunate situation? In the case of $\bmhpc_{m,r}$, it is not clear how to do this in less than $O(m^2)$ bits of communication. One can think of two strategies: either reveal an entire layer, causing the protocol to become aligned at that layer, or solve roughly $\frac{m}{r}$ set intersections instances in each layer, which will cause the protocol to become aligned (follow the pointer of a solved instance) at any given layer with probability $\frac{1}{r}$. We will prove that these two strategies are essentially optimal.

However, in the case of $\bhpc_{m,r}$, it can be seen that $O(\frac{m^2}{r})$ bits of communication suffice. Indeed, $P_C$ and $P_D$ (who speak first) may choose $N = O(\frac{m}{r})$ of their Set-Intersection instances uniformly at random, solve them all by communicating $O(\frac{m^2}{r})$ bits, and then send the solution of all these instances to $P_A$ and $P_B$. Then, in the next round, $P_A$ and $P_B$ begin the standard protocol for solving $\hpc_{m,r}$, meaning, they compute $z_1$ and send it back, then $P_C$ and $P_D$ compute $z_2$ and send it back, etc. If at any point $P_A$ and $P_B$ compute some $z_j$ which is among the $N$ revealed Set-Intersection instances, then they can skip a step, and learn $z_{j+2}$ immediately without communicating. The protocol is now aligned, and can be finished within the remaining number of rounds. The probability that the path taken crosses one of the instances revealed by $P_C$ and $P_D$ is $\Omega(r/ N)$, which can be made arbitrarily close to $1$ with our choice of $N$.

On the other hand, one should not expect to be able to prove any significant lower-bound when $r$ is a sufficiently large fraction of $\sqrt n$. Indeed, in this case, the birthday paradox says that the path being traversed will cycle with high probability. The players may then follow the path (wasting the first round entirely) until they find a cycle, and then they will immediately learn $z_k$.

\medskip\noindent
The intuitive explanation and the upper-bounds given above will lead us to conjecture the following, necessarily tight lower-bounds, which are the main results of this section.

 \begin{theorem} \label{thm:lb-mhpc}
 Any misaligned $r$-round protocol, solving $\bmhpc_{m,r}$ correctly with constant probability requires $\Omega(m^2)$ bits of communication.
 \end{theorem}
 

 \begin{theorem} \label{thm:lb-hpc}
Any misaligned $r$-round protocol, where $r = o(\sqrt m)$, solving $\bhpc_{m,r}$ correctly with constant probability requires $\Omega(\frac{m^2}{r})$ bits of communication.
 \end{theorem}



 A weaker lower-bound of $\frac{m^2}{r^2}$ for HPC was proven in~\cite{AssadiCK19}.
 Our proofs of these tight lower-bounds follow a similar strategy to the earlier
 proof, but with a somewhat modified protocol inspired by~\cite{Yehudayoff20}.
We begin by proving \Cref{thm:lb-hpc} in \Cref{sec:eps-solving-to-mhpc}, because the proof is slightly simpler, and it is enough for our streaming lower-bounds. The proof of \Cref{thm:lb-hpc} appears in \Cref{sec:eps-solving-to-bhpc}.
%
Both proofs follow a similar outline. We first show that a short, misaligned $r$-round protocol for $\bmhpc$ (or $\bhpc$) must be ``almost solving'' an instance of Set-Intersection on one of the rounds, while revealing little information.
%
The notion of ``almost solving'' will be formally defined in \Cref{sec:almost-solving}. 
%
Then, in \Cref{sec:exact-to-eps}, we show that any low-information protocol for ``almost-solving'' Set-Intersection can be used to actually solve Set-Intersection, also with low information. 
%
Putting it all together in \Cref{sec:put-together}, we will find that a misaligned $r$-round protocol for $\bmhpc_{m,r}$ communicating $o(m^2)$ bits (or for $\bhpc_{m,r}$ communicating $o(\frac{m^2}{r})$ bits), would allow us to solve Set-Intersection with $o(m)$ information cost, which is known to be impossible by \cite{JayramKS03}.



\subsection{``Almost Solving''}\label{sec:almost-solving}

The notion of ``almost solving'' is given precisely below.
Intuitively, a protocol $\eps$-solves Set-Intersection if the distribution of
the intersection point $\Set{\restar} = \rA \cap \rB$ is sufficiently altered by
the execution of the protocol.
Meaning, the distribution of $\restar$ conditioned on knowing the protocol (and
possibly one of the inputs) is sufficiently far-away from the distribution of
$\restar$ without this knowledge.
In~\cite{AssadiCK19}, the chosen notion of \textit{sufficiently far} is the
total variation distance.
Here we will choose positive triangular discrimination.

\begin{definition}[$\eps$-solving Set-Intersection, with respect to positive
  triangular discrimination]\label{def:eps-solving}
Let $\cD$ be a distribution on inputs $(A, B)$ for Set-Intersection.
A randomized protocol $\pi$ \textit{internally $\eps$-solves Set-Intersection} (with
respect to positive triangular discrimination) on $\cD$ if \textit{at least one}
of the following hold true (where $\Pi$ is the transcript of $\pi$, including
public randomness):
\[
  \Exp*_{\rPi, \rA}{\PTD{\mu(\restar \mid \rPi, \rA)}{\mu(\restar \mid \rA)}} \geq \eps
  \qquad\text{ or } \qquad
  \Exp*_{\rPi, \rB}{\PTD{\mu(\restar \mid \rPi, \rB)}{\mu(\restar \mid \rB)}} \geq \eps.
\]
A protocol $\pi$ \textit{externally $\eps$-solves Set-Intersectiom} on $\cD$ if
the following holds true:
\[
  \Exp*_{\rPi}{\PTD{\mu(\restar \mid \rPi)}{\mu(\restar )}} \geq \eps.
\]
\end{definition}

\begin{observation}\label{obs:eps-solving-remove-conditioning}
    From the joint convexity of positive triangular discrimination (\Cref{prop:convex}), it follows that for any (possibly correlated) random variables $\rA, \rB, \rC$,
\begin{align*}
    \E_{\rB,\rC}[\PTD{\mu(\rA\mid \rB,\rC)}{\mu(\rA\mid \rC)}] %
    & = \E_{\rB}\left[\sum_{c} \Pr[\rC =c] \PTD{\mu(\rA\mid \rB,\rC=c)}{\mu(\rA\mid \rC=c)}\right] \\
    & \ge \E_{\rB}\left[\Lambda\left(\sum_{c} \Pr[\rC =c]\mu(\rA\mid \rB,\rC=c), \sum_{c} \Pr[\rC =c]\mu(\rA\mid \rC=c)\right)\right]\\
    & = \E_{\rB}[\PTD{\mu(\rA\mid \rB)}{\mu(\rA)}].
\end{align*}
That is, removing a random variable from both conditions can only decrease the $\Lambda$-distance. In particular, if a protocol externally $\eps$-solves Set-Intersection on some distribution $\cD$, it also internally $\eps$-solves Set-Intersection on $\cD$.
\end{observation}

\begin{observation}\label{obs:eps-solving-private-randomness}
    It may be felt that a more natural definition of internally $\eps$-solving
    should include any eventual private randomness of the protocol $\pi$.
    However, noting that $\{\restar\} = \rA \cap \rB$ is a function of $\rA$ and
    $\rB$, and that $\rA$ and $\rB$ are independent given $\rPi$ (by
    \Cref{fact:rectangle})
    it follows that the distribution of $\restar$ is independent of the private randomness of $\pi$, when given $\rA$ and $\rPi$. So
    \[
    \Exp*_{\rPi, \rA}{\PTD{\mu(\restar \mid \rPi, \rA)}{\mu(\restar \mid \rA)}} = \Exp*_{\rPi, \rA}{\PTD{\mu(\restar \mid \rPi, \rR_A, \rA)}{\mu(\restar \mid \rA)}}
    \]
    where $\rR_A$ is Alice's randomness, and likewise on Bob's side.
    That is, including the private randomness gives the same definition.
\end{observation}


\subsection{Hard Distributions} \label{sec:distribution}


Next, we define two hard distributions for $\cD_\bmhpc$ and $\cD_\bhpc$,  over inputs to $\bmhpc$ and $\bhpc$, respectively. 


\begin{tbox}
    \textbf{Distribution $\cD_\bmhpc$} on tuples $(\rA,\rB,\rC,\rD)$ from the universes $\cX$ and $\cY$: 

    \begin{enumerate}
		\item For any $i\in [r]$, and any $x \in \cX$, sample $(A^i_x,B^i_x) \sim \cD_\SI$ from the universe $\cY$ \emph{independently}. 
		\item For any $i \in [r]$, and any $y \in \cY$, sample $(C^i_y,D^i_y) \sim \cD_\SI$ from the universe $\cX$ \emph{independently}. 
    \end{enumerate}
\end{tbox}

\begin{tbox}
    \textbf{Distribution $\cD_\bhpc$} \cite{AssadiCK19} on tuples $(\rA,\rB,\rC,\rD)$ from the universes $\cX$ and $\cY$: 

    \begin{enumerate}
		\item For any $x \in \cX$, sample $(A_x,B_x) \sim \cD_\SI$ from the universe $\cY$ \emph{independently}. 
		\item For any $y \in \cY$, sample $(C_y,D_y) \sim \cD_\SI$ from the universe $\cX$ \emph{independently}. 
    \end{enumerate}
\end{tbox}

It was shown in~\cite{AssadiCK19} that misaligned randomized $r$-round protocols
for $\hpc_{m,r}$ need $\frac{m^2}{r^2}$ bits of communication under $\cD_\bhpc$.
In \Cref{sec:eps-solving-to-bhpc}, we improve this result and show that
misaligned randomized $r$-round protocols for $\bhpc_{m,r}$ (an easier problem
than $\hpc_{m,r}$) need $\frac{m^2}{r}$ bits of communication.


We restate the following simple observations from~\cite{AssadiCK19} about
$\cD_\bhpc$.
Analogous statements can be made for $\cD_\bmhpc$.

\begin{observation}[\cite{AssadiCK19}]\label{obs:distHPC}
	Distribution $\cD_\bhpc$ is \underline{not} a product distribution. However, in this distribution: 
	\begin{enumerate}[label=(\roman*)]
		\item The inputs to $P_A$ and $P_B$ are {independent} of the inputs to
      $P_C$ and $P_D$, i.e., $(\rA,\rB) \perp (\rC,\rD)$. 
		\item For any $x \in \cX$, $(\rA_x,\rB_x)$ is independent of all other
      $(\rA_{x'},\rB_{x'})$ for $x'\neq x \in \cX$.
      Similarly for all $y,y' \in \cY$ and $(\rC_y,\rD_y)$ and $(\rC_{y'},\rD_{y'})$. 
	\end{enumerate}
\end{observation}

Based on this observation, we also have the following simple property. 

\begin{proposition}[\cite{AssadiCK19}]\label{prop:hpc-rectangle}
  Let $\prot_\bhpc$ be any deterministic protocol for $\bhpc_{m,r}$ on
  $\cD_\bhpc$.
  Then, for any transcript $\Prot$ of $\prot_\bhpc$, $(\rA,\rB)
  \perp (\rC,\rD) \mid \rPi = \Prot$. 
\end{proposition}
\begin{proof}
  Follows from the rectangle property of the protocol $\prot_\bhpc$
  (Fact~\ref{fact:rectangle}).
  In particular, the same exact argument as in the two-player case implies that
  if $[(A, B),(C, D)]$ and $[(A', B'),(C', D')]$ are mapped to the same
  transcript $\Prot$, then $[(A, B),(C', D')]$ and $[(A', B'),(C, D)]$ are
  mapped to $\Prot$ as well.
  Hence, since $(\rA, \rB) \perp (\rC,\rD)$ by \Cref{obs:distHPC}, the
  inputs corresponding to the same protocol would also be independent of each
  other, namely, $(\rA, \rB) \perp (\rC, \rD) \mid \rPi = \Prot$. 
\end{proof}







\subsection{\texorpdfstring{Short Protocols for $\bmhpc$ Give Low-Information $\eps$-Solvers for $\SI$}{Short Protocols for BMHCP Give Low-Information eps-Solvers for Set Intersection}}\label{sec:eps-solving-to-mhpc}

We now show that if we have a short, $r$-round protocol for $\bmhpc_{m,r}$ on $\cD_\bmhpc$, then during one of its rounds it must be externally (and hence also internally) $\eps$-solving Set-Intersection on $\cD_\SI$, while revealing little information. This is proven by making precise the following intuitive line of thinking, the statement of which is formally captured by \Cref{lem:mhpc-dist}. 

\begin{itemize}
  \item
      By the end of round $j$, the speaking pair of players would like to
      have learned $\{\rZ_{j} \} = \rA^j_{\rZ_{j-1}} \cap \rB^j_{\rZ_{j-1}}$
      (for odd $j$, or $\{\rZ_{j} \} = \rC^j_{\rZ_{j-1}} \cap \rD^j_{\rZ_{j-1}}$,
      for even $j$).
      Now suppose that in this round they have very little information about
      $\rZ_{j-1}$.
      Let us say, for example, that the distribution of
      $\rZ_{j-1} \mid \rE_{j-1}, \rPi_{j-1}$ is close to uniform.
      This certainly holds when $j = 2$, since
      $\mu(\rZ_1 | \rE_1, \rPi_1) = \mu(\rZ_1)$ is
      completely uniform (because the wrong pair of players speak first and the
      way we construct $\cD_\bmhpc$).
    


  \item Now suppose that despite this, somehow, the two speaking players
    nonetheless manage to learn a lot of information about $\rZ_{j}$. Then they
    have ``almost solved'' an \textit{unknown} instance (an instance belonging
    to the other players, no less), out of the possible remaining instances of
    Set-Intersection.
    Because they have little information about $\rZ_{j-1}$,
    meaning $\rZ_{j-1}$ is close to uniform, they have thus ``almost solved'' an
    unknown, nearly-uniformly-chosen instance. 

  \item Then, intuitively, they must have ``almost solved'' this unknown instance of Set-Intersection without revealing a lot of information about it. This turns out to be impossible, as we will see in \Cref{sec:exact-to-eps}. 
\end{itemize}


One difficulty in implementing this intuitive strategy is in finding the right measure of information that allows for one to carry through the entire argument. In \cite{AssadiCK19}, \textit{total variation distance from uniform} was used as a measure of how much the speaking players know about which Set-Intersection instance they are trying to solve. But TVD incurs the quadratic loss inherent in Pinsker's inequality, leading to a lower-bound of $\frac{m^2}{r^2}$. The same quadratic loss appeared originally in Nisan and Wigderson's lower-bounds for pointer chasing \cite{NisanW93}, and this led to a three-decade-long gap between known upper and lower-bounds for pointer chasing. This gap was finally closed by Amir Yehudayoff \cite{Yehudayoff20}, using positive triangular discrimination. This suggested the possibility of using Yehudayoff's approach to improve the lower-bound of \cite{AssadiCK19}, and it is this plan that we realize here.

\bigskip\noindent
The following lemma states that no low-information protocol (in the sense of information cost) for $\bmhpc$ can reveal a lot of information (in the sense of $\Lambda$-distance-from-uniform) about a uniformly chosen $\setint$ instance at any layer $j$. The proof shows that any such protocol would give us a low-information protocol for $\eps$-solving $\setint$, which we will show is impossible in \Cref{sec:exact-to-eps}.

\begin{restatable}%[Progress Lemma]
{lemma}{progress-multi} \label{lem:progress}
There is a constant $c > 0$ such that, for any protocol for $\bmhpc_{m,r}$
running on input $(\rA,\rB,\rC,\rD)$ sampled from $\cD_{\bmhpc}$, and for any
$j \in [r]$, we have:
\begin{align*}
\E_{(\rE_j, \rPi_j)}\E_{x \sim \cU_\cX}\left[\Lambda(\mu(\rT^j_x \mid \rE_j,
\rPi_j), \mu(\rT^j_x))\right] 
%
\leq  
%
c \cdot \frac{\mi{\rPi}{\rA^j \mid \rA^{< j}, \rB} + \mi{\rPi}{\rB^j \mid \rB^{> j}, \rA}}{m^2}\\
\E_{(\rE_j, \rPi_j)}\E_{y \sim \cU_\cY}\left[\Lambda(\mu(\rT^j_y \mid \rE_j,
\rPi_j), \mu(\rT^j_y))\right] 
%
\leq  
%
c \cdot \frac{\mi{\rPi}{\rC^j \mid \rC^{< j}, \rD} + \mi{\rPi}{\rD^j \mid \rD^{> j}, \rC}}{m^2}\\
\end{align*}
\end{restatable}

\begin{proof}
    Take any given $j \in [r]$, and consider the following protocol for
    $\eps$-solving $\setint$ on $\cD_{\SI}$:
    
	\begin{abox}
    \textbf{Protocol $\protSI$}: The protocol for $\eps$-solving $\SI$ using a protocol $\protHPC$ for $\bmhpc_{m,r}$. 
    
     \smallskip
    
    \textbf{Input:} An instance $(X, Y) \sim \cD_\SI$ with $X, Y\subseteq\cY$.
    
   % \algline
    
    \begin{enumerate}
    	\item \textbf{Sampling the instance.} Alice and Bob create an instance
        $(A, B, C, D)$ of $\bmhpc_{m,r}$ as follows: %The players create an instance $(A,B)$ of $\SI$ as follows:
	\begin{enumerate}
	
 \item Using \underline{public coins}, Alice and Bob sample an index
   $I \in [m]$ uniformly at random, and  
	Alice sets $A^j_{x_I} = X$ and Bob sets $B^j_{x_I} = Y$ using their given inputs in $\SI$. 
	
 \item Using \underline{public coins}, Alice and Bob sample $A^j_{x_1}, \ldots, A^j_{x_{I-1}}$ from the A-side marginal of $\cD_\SI$, and they sample $B^j_{x_{I+1}}, \ldots, B^j_{x_{m}}$ from the B-side marginal of $\cD_\SI$. 
	
 \item Using \underline{private coins}, Alice samples $A^j_{x_r}$ for $r \in \{I+1, \ldots, n\}$ so that $(A^j_{x_r},B^j_{x_r}) \sim \cD_\SI$. Similarly Bob samples $B^j_{x_\ell}$ for $\ell \in \{1, \ldots, I-1\}$. This completes the construction of $(A^j, B^j)$. 
        
 \item Using \underline{public coins}, Alice and Bob sample $A^{< j}, B^{> j}, C$ and $D$ from their respective marginals in $\cD_\bmhpc$ (this is possible by independence).
 \item Using \underline{private coins}, Alice samples $A^{> j}$ conditioned on the publicly sampled $B^{>j}$, and Bob samples $B^{< j}$ conditioned on the publicly sampled $A^{<j}$. 

        \end{enumerate}
  %      (It should be evident that the distribution of instances $(\bA,\bB,\bC,\bD)$ created by $\protSI$ is $\distHPC$.) 
        
        
	\item The players now run the first $j$ rounds of the protocol $\prot_\bmhpc$ on the instance $(A,B,C,D)$, by Alice playing $P_A$, Bob playing $P_B$, and both Alice and Bob simulating $P_C$ and $P_D$ with no communication (this is possible as both Alice and Bob know $(C, D)$ entirely). 
 
    Let us write $\Pi_\SI$ to denote the messages of $P_A$ and $P_B$, which were actually sent by Alice and Bob, and $\Pi_{\le j}$ to denote the entire simulated transcript of $\prot_\bmhpc$, including the messages of $P_C$ and $P_D$ (which Alice and Bob learned without communicating).
    \end{enumerate}
\end{abox}

From a simple direct-sum argument, we get an upper-bound on the information cost of $\pi_\SI$:

\begin{claim} \label{clm:direct-sum}
    \[ \ICost{\pi_\SI}{\dsi} = \frac{\mi{\rPi_{\le j}}{\rA^j \mid \rA^{< j}, \rB}
    + \mi{\rPi}{\rB^j \mid \rB^{> j}, \rA}}{m}. \]
\end{claim}

\begin{proof}
    The internal information cost of $\prot_\SI$ is
    \[
      \IC{\prot_\SI}{\cD_\SI}
      := \mi{\rX}{\rProtSI \mid \rY, \rR, \rR_B}
      + \mi{\rY}{\rProtSI \mid \rX, \rR, \rR_A},
    \]

    where $R, R_A, R_B$ are the public, Alice's, and Bob's randomness.
    Let us bound the first term, as the second term can be bounded in just the
    same way.
    We then derive:   
\begin{align*}
	\mi{\rPi_\SI}{\rX \mid \rY, \rR, \rR_B} %
    & = \mi{\rPi_{\le j}}{\rX \mid \rY, \rI,\rA^j_{<\rI},\rB^j_{>\rI},\rA^{< j},\rB^{> j},\rC,\rD, \rB^j_{< \rI}, \rB^{< j}}\\
    & = \mi{\rPi_{\le j}}{\rA^j_\rI \mid \rI,\rA^j_{<\rI},\rA^{< j},\rB,\rC,\rD}\\
    & =  \sum_{i=1}^m \frac{1}{m} \mi{\rPi_{\le j}}{\rA^j_i \mid \rI=i,\rA^j_{<i},\rA^{< j},\rB,\rC,\rD}\\
    & =  \sum_{i=1}^m \frac{1}{m} \mi{\rPi_{\le j}}{\rA^j_i \mid \rA^j_{<i},\rA^{< j},\rB,\rC,\rD}\\
    & =  \frac{1}{m} \mi{\rPi_{\le j}}{\rA^j \mid \rA^{< j},\rB,\rC,\rD}\\
    & =  \frac{1}{m} \mi{\rPi_{\le j}}{\rA^j \mid \rA^{< j},\rB}. 
\end{align*}
All equalities are by definition, except for the fourth which is because none of
the remaining distributions depend on the choice of $\rI$, and the last, which
is because $\rC, \rD$ are independent of $\rA^j$, even when conditioned on
$\rPi_{\le j}, \rA^{<j}, \rB$.
\end{proof}

On the other hand, we also get an $\eps$-solver, for $\eps$ being exactly the $\Lambda$ distance-from-uniform of a uniformly chosen coordinate in the $j$-th layer:

\begin{claim} \label{clm:average-coordinate}
$\pi_\SI$ internally $\eps$-solves Set-Intersection on $\dsi$ where:
\[
  \eps \geq \E_{(\rE_j, \rPi_j)}
  \Exp*_{x \sim \cU_\cX} {\PTD{\mu(\rT^j_x \mid \rE_j, \rPi_j)}{\mu(\rT^j_x)}}.
\]
\end{claim}


\begin{proof} 
     By \Cref{def:eps-solving} and \Cref{obs:eps-solving-private-randomness}, we have an $\eps$-solver on Bob's side, for
     \[
       \eps = \E[\Lambda(\mu(\rT \mid \rY, \rPi_\SI,
     \rI, \rR, \rR_B), \mu(\rT \mid \rY, \rI, \rR, \rR_B))], 
     \]
     where $T = X \cap Y = A^j_I \cap B^j_I$,
     $I$ and $R = A^j_{<I}, B^j_{>I}, A^{< j}, B^{> j}, C, D$ are the public
     randomness, and
     $R_B = B^j_{< I}, B^{< j}$ is Bob's private randomness. 
     This is exactly equal to
     \[
     \E[\Lambda(\mu(\rT \mid \rY, \rE_j, \rPi_{j}, \rI, \rR, \rR_B),
     \mu(\rT \mid \rY, \rI, \rR, \rR_B))], \]
     because $E_j, \Pi_j$ is a function of $\Pi_\SI, I, R, R_B$ (Bob knows both
     $A^{<j}$, $B^{<j}$ and $C$, $D$, so he also knows $Z_{< j}$, and
     $\Pi_{\le j}$).
     In turn, this equals
     \begin{align*}
       \E_{\rE_j,\rPi_j, i \sim \uni_{\brac{m}}}
       \Exp*{\PTD{\mu(\rT^j_{x_i} \mid \rE_j, \rPi_{j}, \rI = i, \rY, \rR,\rR_B)}
       {\mu(\rT^j_{x_i}\mid \rY, \rI=i,\rR, \rR_B)}}
     \end{align*}
     Observe that none of the distributions above depend on the event $\rI = i$. So this quantity is equal to
     \begin{align*}
       \E_{\rE_j,\rPi_j, i \sim \uni_{\brac{m}}}
       \Exp*{\PTD{\mu(\rT^j_{x_i} \mid \rE_j, \rPi_{j}, \rY, \rR,\rR_B)}
       {\mu(\rT^j_{x_i}\mid \rY, \rR, \rR_B)}}
     \end{align*}
     Now, by \Cref{obs:eps-solving-remove-conditioning}, this is greater or equal to
     \[
       \E_{\rE_j,\rPi_j, i \sim \uni_{\brac{m}}}
       \Exp*{\PTD{\mu(\rT^j_{x_i} \mid \rE_j, \rPi_{j} )}{\mu(\rT^j_{x_i})}}. \qedhere
     \]
\end{proof}


From \Cref{cor:si-ic-lb-external} proven in the next section, it now follows:
  \begin{align*}
    \frac{\mi{\rPi_{\le j}}{\rA^j \mid \rA^{< j}, \rB} + \mi{\rPi_{\leq j}}{\rB^j \mid \rB^{> j}, \rA}}{m} & \ge c \cdot \eps m\\
    & \geq c m \cdot \E_{(\rE_j, \rPi_j)}\E_{x \sim \cU_\cX}
    \left[\Lambda(\mu(\rT^j_x \mid \rE_j, \rPi_j), \mu(\rT^j_x)) \right].
  \end{align*}
  This concludes the proof of \Cref{lem:progress}.
\end{proof}



\bigskip\noindent
It is now possible to prove that there is no low-information protocol for solving $\bmhpc$. The argument was roughly outlined above: if we assume by induction that the $\Lambda$ distance-from-uniform is small for the answer at the previous layer, then we are essentially trying to solve a uniformly random instance at the current layer. And we just proved that this is not possible.

\begin{lemma}\label{lem:mhpc-dist}
For the same constant $c > 0$ of the \Cref{lem:progress}, and for any protocol
for $\bmhpc_{m,r}$ running on input $(A, B, C, D)$ sampled from
$\cD_{\bmhpc}$, the following statement holds for all $j \in
[r]$:%
\footnote{Intuitively, the statement says that, if $\pi_\bmhpc$
communicates little, then the distribution of $Z_j$, conditioned on the
information available at the end of round $j$, is close to uniform.}
\[
  \E_{(\rE_j, \rPi_j)}\left[\Lambda(\mu(\rZ_j \mid \rE_j, \rPi_j), \mu(\rZ_j))\right]
  \leq 6 c \cdot \frac{\mi{\rPi}{\rA^{\le j},\rC^{\le j} \mid \rB,\rD} + \mi{\rPi}{\rB^{\le j},\rD^{\le j} \mid \rB^{> j}, \rD^{> j}, \rA,\rC}}{m^2}.
\] 

\end{lemma}





This is proven by induction. The base case is easy to see: note that $\mu(\rZ_1
\mid \rE_1, \rPi_1) = \mu(\rZ_1)$ as $\rZ_1$ is independent of $\rE_1 = z_0$
which is fixed and $\rZ_1$ is independent of $\rPi_1$ as $\Pi_1$ is a function
of $(C,D)$.
We obtain case $j$ from case $j - 1$ in two steps: we first show that if the
protocol reveals a lot of information about the variable $\rZ_{j}$, then it must
essentially be $\eps$-solving a \textit{uniformly random instance}. This is
because $Z_{j} = T_{Z_{j-1}}$ and $\rZ_{j-1}$ is, under the inductive case
$j-1$, close to uniform.
We then appeal to the previous lemma.

\begin{proof}
The information---measured by $\Lambda$-distance to uniform---that the $j$-th
round of the protocol reveals about $\rZ_j = \rT_{\rZ_{j-1}}$ is the
information revealed about an instance sampled according to $\rZ_{j-1}$:
  \begin{align*}
		\E_{(\rE_j,\rProt_j)} \Bracket{\Lambda(\mu(\rZ_j \mid
    \rE_j,\rProt_j),\mu(\rZ_j))} &=\E_{(\rE_j,\rProt_j)}
    \Bracket{\Lambda(\mu(\rZ_j \mid \rE_j),\mu(\rZ_j))}\\
		&\hspace{-1.5cm}= \E_{(\rZ_{<j},\rProt_{<j})} \Bracket{\Lambda(\mu(\rZ_j
    \mid \rZ_{<j},\rProt_{<j}),\mu(\rZ_j))} \tag{by definition of
  $\rE_j := (\rZ_{<j},\rProt_{<j})$, and since $\rZ_j$ and $\rPi_j$ are
independent given $\rE_j$.} \\
		&\hspace{-1.5cm}=\E_{(\rZ_{<j},\rProt_{<j})}
    \Bracket{\Lambda(\mu(\rT^j_{\rZ_{j-1}} \mid
    \rZ_{<j-1},\rZ_{j-1},\rProt_{<j}),\mu(\rZ_j))}\tag{by definition, the pointer $Z_j = T_{Z_{j-1}}$}\\
		&\hspace{-1.5cm}= \E_{(\rZ_{<j-1},\rProt_{<j})} \; \E_{z \sim \rZ_{j-1} \mid
    \paren{\rZ_{<j-1},\rProt_{<j}}}\; \Bracket{\Lambda(\mu(\rT^j_{z} \mid
  \rZ_{<j-1},\rProt_{<j}),\mu(\rZ_j))}.
	\end{align*}
  We removed the conditioning on $\rZ_{j-1} = z$ in the first distribution
  because, conditioned on knowing $\rZ_{<j-1},\rProt_{<j}$, for any fixed value of
  $z$ the value of $\rT_{z}$ is independent of the value of $\rZ_{j-1}$:
  \begin{itemize}
	    \item If $j-1$ is odd, $\rT_{z}$ is a function of  $(\rC, \rD)$ and
        $\rZ_{j-1}$ is a function of $(\rA, \rB)$.
            \item If $j-1$ is even, $\rT_{z}$ is a function of $(\rA,\rB)$ and
              $\rZ_{j-1}$ is a function of $(\rC, \rD)$.
            \item And, furthermore, $(\rA,\rB)$ and $(\rB,\rD)$ are independent
              given $\rZ_{<j-1}, \rPi_{<j}$. 
	\end{itemize}

    
	\medskip\noindent
    \Cref{lem:progress} says that if the protocol were to reveal information
    about the solution $\rT_z$ of a truly uniform instance $z$, then this would
    give us a protocol for $\eps$-solving set disjointness. But, by induction,
    we know that the chosen instance $z \sim \rZ_{j-1}$ is close to uniform. It
    should then follow from \Cref{itm:lambda-loss} that if the protocol reveals
    some information about this almost uniform instance, it must also reveal
    information about a truly uniform instance. Let us now make this precise.
    
  Without loss of generality suppose that $j-1$ is odd and hence $z \in \cY$.
  Then we can upper bound the last expression above
  by:\footnote{\label{fn:tvd}In the first inequality it should be understood how
  $\Lambda$ is playing two different roles. On the one hand it is being used as
an information-theoretic quantity (the distance to uniform), and on the other it
is being used in a capacity similar to how total variation distance appears in
the obvious upper-bound $\E_\mu[f(x)] \le \E_\nu[f(x)] + \|\mu - \nu\|_1 \cdot
\max f(x)$.}

	\begin{align*}
          &\E_{(\rE_j,\rPi_j)} \Bracket{\Lambda(\mu(\rZ_j \mid \rE_j,\rPi_j),\mu(\rZ_j))}  \\
          % 
          &\qquad \qquad \leq \E_{(\rZ_{<j-1},\rPi_{<j})} \bracket{6 \cdot \E_{\paren{z \sim \cU_{\cY}}}\Bracket{\Lambda(\mu(\rT^j_{z} \mid \rZ_{<j-1},\rPi_{<j}),\mu(\rZ_j))}} \\
          % 
          &\qquad \qquad  \qquad \qquad \qquad  \qquad + \E_{(\rZ_{<j-1},\rPi_{<j})} \Bracket{\Lambda(\mu(\rZ_{j-1} \mid \rE_{j-1},\rPi_{j-1}),\cU_\cY)} %
            \tag{by \Cref{itm:lambda-loss} and linearity of expectatiom}
          \\
          % 
          &\qquad \qquad = 6 \cdot\E_{(\rZ_{<j-1},\rPi_{<j})} \bracket{ \E_{\paren{z \sim \cU_{\cY}}}\Bracket{\Lambda(\mu(\rT^j_{z} \mid \rZ_{<j-1},\rPi_{<j}),\mu(\rZ_j))}} \\
          &\qquad \qquad  \qquad \qquad \qquad  \qquad + \E_{(\rZ_{<j-1},\rPi_{<j})} \Bracket{\Lambda(\mu(\rZ_{j-1} \mid \rE_{j-1},\rPi_{j-1}),\mu(\rZ_{j-1}))}\tag{by definition $z_{j-1} \in \cY$ and $\mu(\rZ_{j-1}) = \cU_\cY$.}\\
          % 
          &\qquad \qquad\leq \bigg[6 \cdot \underbrace{c \cdot \frac{\mi{\rPi}{\rA^j \mid \rA^{< j}, \rB} + \mi{\rPi}{\rB^j \mid \rB^{> j}, \rA}}{m^2}}_{\text{\Cref{lem:progress}}} \\
          &\qquad \qquad\qquad 
            + \underbrace{6 c \cdot \frac{\mi{\rPi}{\rA^{\le j-1},\rC^{\le j-1} \mid \rB,\rD}
           + \mi{\rPi}{\rB^{\le j-1},\rD^{\le j-1} \mid \rB^{> j-1}, \rD^{> j-1}, \rA,\rC}}{m^2}}_{\text{induction hypothesis (case $j-1$)}}\bigg]\\
    % 
    % 
          & \qquad \qquad= \bigg[6c \cdot \frac{\mi{\rPi}{\rA^j \mid \rA^{< j},\rC^{< j}, \rB,\rD} + \mi{\rPi}{\rB^j \mid \rB^{> j},\rD^{>j}, \rA,\rC}}{m^2} \\
          &\qquad \qquad\qquad 
            + 6 c \cdot \frac{\mi{\rPi}{\rA^{\le j-1},\rC^{\le j-1} \mid \rB,\rD}
            + \mi{\rPi}{\rB^{\le j-1},\rD^{\le j-1} \mid \rB^{> j-1}, \rD^{> j-1}, \rA,\rC}}{m^2}\bigg]\\
          &         \qquad \qquad
          \leq
          \bigg[6 c \cdot \frac{\mi{\rPi}{\rA^j,\rC^j\mid \rA^{< j},\rC^{< j}, \rB,\rD} + \mi{\rPi}{\rB^j,\rD^j \mid \rB^{> j},\rD^{>j}, \rA,\rC}}{m^2} \\
          &\qquad \qquad\qquad 
            + 6 c \cdot \frac{\mi{\rPi}{\rA^{\le j-1},\rC^{\le j-1} \mid \rB,\rD}
            + \mi{\rPi}{\rB^{\le j-1},\rD^{\le j-1} \mid \rB^{> j-1}, \rD^{> j-1}, \rA,\rC}}{m^2}\bigg]\\
          & \qquad\qquad =
            6 c \cdot \frac{\mi{\rPi}{\rA^{\le j},\rC^{\le j} \mid \rB,\rD}
           + \mi{\rPi}{\rB^{\le j},\rD^{\le j} \mid \rB^{> j}, \rD^{> j},
         \rA,\rC}}{m^2},
	\end{align*}
where the last inequality follows from the data processing inequality
(\Cref{fact:it-facts}~\Cref{part:data-processing}).
This concludes the proof.
\end{proof}




\subsection{\texorpdfstring{$\eps$-Solvers for $\SI$ Give Exact Solvers}{eps-Solvers for Set Intersection Give Exact Solvers}}\label{sec:exact-to-eps}



We now show a lower-bound against protocols trying to $\eps$-solve
Set-Intersection, by reducing from exactly solving Set-Intersection:


\begin{theorem}\label{thm:si-ic-lb} For any constant $\gamma \in (0,1)$, there exists a constant $c \ge 1$ such that, given a protocol $\pisi$ which internally 
$\eps$-solves Set-Intersection (in $\ptd$ distance) on $\dsi$ with
$\eps \ge \frac{8}{m}$,
we may obtain a protocol $\pi$ that exactly solves Set-Intersection on $\dsi$,
such that
$\ICost{\pi}{\dsi} \leq {\frac{c}{\eps} \cdot \ICost{\pisi}{\dsi} + \gamma m}$.
\end{theorem}

From the known fact (\Cref{prop:SI-hard}) that Set-Intersection has information
complexity $\Omega(m)$ on $\dsi$ we get the following:

\begin{corollary}\label{cor:si-ic-lb-external}
  There exists a constant $c > 0$ such that
  any protocol $\pisi$ that internally $(\ptd, \eps)$-solves $\setint_m$ on
  inputs from the distribution $\dsi$, with $\eps \geq 8 / m$, has $\ICost{\pisi}{\dsi} \geq c \cdot \eps m$.
\end{corollary}

An analogue of \Cref{thm:si-ic-lb} was proven in
\cite{AssadiR20}, where $\eps$-solving was with respect to total variation
distance instead of triangular discrimination.
It was then shown that any protocol for $\eps$-solving Set-Intersection in this
TVD sense gives a protocol for Set-Intersection with a $O(1/\eps^2)$
blowup in information cost.
\Cref{thm:si-ic-lb}, instead, gives a blowup of $O(1/\eps)$.
So our assumption is stronger, but our conclusion is also stronger.

\begin{proof}[Proof of \Cref{thm:si-ic-lb}]
We begin by recalling the notation we use for the Set-Intersection problem:
Alice and Bob's input sets are
denoted by $X$ and $Y$ respectively, and their corresponding random variables
by $\rX$ and $\rY$.
The unique element in the intersection $X \cap Y$ is denoted by $\estar$,
and its random variable is $\restar$.
Throughout this proof, for a permutation $\sigma$ and a set $S$,
$\sigma(S)$ denotes the set $\Set{ \sigma(s) \mid s \in S }$.

\noindent
Without loss of generality, we may assume that we have a protocol $\pisi$ that
internally $\eps$-solves from \emph{Alice}'s side
Set-Intersection on inputs distributed according to $\dsi$, that is,
\[
  \rPTD = 
  \Exp*_{\rX, \rPi}{
     \PTD*{ \Dist{ \restar \given \rPi, \rX}}
     {\Dist{ \restar, \rX}}
  } \geq \eps,
\]
where $\eps \geq 8 / m$.
We will use $\pisi$ to design a protocol $\pi$ which solves Set-Intersection
on $\dsi$.

For now, suppose that Alice and Bob proceed as follows: they sample a uniformly
random permutation $\Perm$ of $[m]$ using public randomness, and run the
$\eps$-solving protocol $\pisi$ on $(\Perm(X), \Perm(Y))$, obtaining a
transcript $\Pi$.
Alice now knows the induced distribution of the intersection
point $\Dist{\rPerm(\estar) \given \rPi = \Pi, \rPerm(X) = \Perm(X)}$, conditioned
on the transcript she has seen, and Bob knows the analogous distribution with
his input fixed.
Without the conditioning on the transcript, we know by symmetry of $\dsi$ that
$\Dist{\rPerm(\estar) \given \rPerm(X) = \Perm(X)}$ is uniform over $\Perm(X)$.
In fact, we have the following observation about the distribution of inputs
generated by applying $\rPerm$.

\begin{observation}\label{obsv:dsi}
  For any fixed $X$ and $Y$ such that $\card{X} = \card{Y} = m / 4$, and
  $\card{X \cap Y} = 1$,
  and a uniformly random permutation $\rPerm$,
  the instance $(\rPerm(X), \rPerm(Y))$ is distributed according to $\dsi$.
\end{observation}

To see this, consider the fixed values $X = \Set{1, \ldots m / 4}$ and
$Y = \Set{m / 4, \ldots m / 2 - 1}$.
Because $\rPerm$ is a uniformly random permutation, it maps $m / 4$ to a
uniformly random element of $\brac{m}$, and the rest of $X$ and $Y$ to
uniformly random subsets of size $m / 4 - 1$.
By symmetry, this argument applies to any fixed $X$ and $Y$ of size $m / 4$
intersecting in exactly one element.


\subsubsection{A single round of \texorpdfstring{$\pi$}{the protocol}}\label{subsec:single-round}

We start by describing a single round of the protocol $\pi$.
On the input $(X, Y)$:
\begin{enumerate}
  \item Alice and Bob publicly sample a random permutation $\Perm$ of
    $\bracket{m}$, and apply it to $X$ and $Y$ to obtain $\Perm(X)$ and
    $\Perm(Y)$ respectively.
  \item They run $\pisi$ on $\Perm(X)$ and $\Perm(Y)$ to obtain a transcript
    $\Pi$.
    By doing so, Alice learns a distribution $\qpi_1, \ldots , \qpi_m$, where
    \[ \qpi_i = \Prob{ \rPerm(\estar) = i \given \rPi = \Pi, \rPerm(X) = \Perm(X)} . \]
  \item Alice uses $\qpi$ to assign ``scores'' to each element $e$ of
    $\brac{m}$, which reflect the likelihood that $\Set{e}$ is the
    intersection of $X$ and $Y$.
\end{enumerate}

This strategy follows the general idea of~\cite{AssadiR20}, but where they assigned the score $1$ to all $e$ such that $\qpi_{\Perm(e)}$ is in the largest-probability half of $\qpi_1, \ldots , \qpi_m$ (and this works for $\eps$-solving under TVD), we assign scores in a smoother fashion.
Let $i = \Perm(e)$ and $n = m / 4$ for brevity.
Then if $\qpi_{i} \leq 1/n$, we give $e$ a score of $0$.
Otherwise, its score is $(\qpi_i - 1/n) / (\qpi_i + 1/n)$.
This choice of the ``scoring function'' may seem odd at this juncture, so
we offer two comments about it:
\begin{itemize}
  \item Note that this is equal to $1 - (2 / n) / (\qpi_i + 1/n)$, which is an
    increasing function of $\qpi_i$, and hence passes a basic sanity test
    (i.e.\ if $\pisi$ says an element is more likely to be $\Perm(\estar)$,
    we believe it).
  \item 
    The scoring function was chosen to make the expected score of $\estar$
    larger than the expected score of any non-$\estar$ element by at least
    $\rPTD$, and the variance of every element's score turned out to be small.
    %It is a mystery to us why the math works out.
    %\parth{I want to leave this last line in.}
\end{itemize}

We package the properties of this scoring function into a few claims below.
Throughout these claims, we will use the shorthand notation
$\rA \mid B$ to mean $\rA \mid \rB = B$, where $\rA$ and $\rB$ are random
variables, and $B$ is some fixed value that $\rB$ can take.

\begin{claim}\label{clm:exp-score-diff}
  For any non-$\estar$ element $e \in \brac{m}$,
  \[ \Exp_{\rPerm} {\score(\estar) - \score(e)} \geq \rPTD / 2. \]
\end{claim}

\begin{proof}
  The proof is in two parts; first we bound $\score(\estar)$ under some
  appropriate conditioning, and then $\score(e)$ for non-$\estar$ elements.

  \noindent
  Fix any $(X,Y)$ with $|X| = |Y| = n$ and $|X\cap Y| = 1$.
  Consider a fixed choice of $\Perm(X)$ and the transcript $\Pi$, and hence
  a distribution $\qpi$.
  Let $P$ denote the set of indices where $\qpi_i > 1/n$ (i.e.\ those that lead
  to positive scores).
  Note that $P \subseteq \Perm(X)$; because we condition on $\Perm(X)$
  (and more) to obtain $\qpi$.
  By definition, the probability that $\estar$ is mapped to $i$ is $q_i$, and
  hence:
  \[
    \Exp*_{\rPerm \given \Perm(X), \Pi} {\score(\estar)}
    = \sum_{i \in P} \qpi_i \cdot \frac{\qpi_i - 1/n}{\qpi_i + 1/n}.
  \]
  On the other hand, fixing both $\Perm(X)$, $\Perm(Y)$ (and hence
  also $\Perm(\estar)$) and $\Pi$, fixes the distribution $\qpi$,
  but leaves completely free the map
  $\Perm: X \setminus \Set{\estar} \to \Perm(X) \setminus \Set{\Perm(\estar)}$.
  So we find that:
 
    $\rPerm(e)$ is uniformly random over
      $\Perm(X) \setminus \Set{\Perm(\estar)}$ for any non-$\estar$ element
      $e \in X$.
  
  
  This means that for any $i \in \brac{m}$,
  \[
    \Prob_{ \rPerm \given \Perm(X), \Perm(Y), \Pi } {\rPerm(e) = i} \leq
    \frac 1{n - 1}.
  \]
  Hence we can compute the expected score of $e$ as:
  \begin{align*}
    \Exp_{\rPerm \given \Perm(X), \Perm(Y), \Pi}{\score(e)}
    &= \sum_{i \in P}  \Prob{\rPerm(e) = i} \cdot \frac{\qpi_i - 1/n}{\qpi_i + 1/n}
    \leq \sum_{i \in P}  \frac 1{n - 1} \cdot \frac{\qpi_i - 1/n}{\qpi_i + 1/n}\\
    &= \sum_{i \in P}  \pparen*{\frac 1{n} + \frac 1{n(n - 1)}} \cdot \frac{\qpi_i - 1/n}{\qpi_i + 1/n}
    \leq \frac 1n + \sum_{i \in P}  \frac 1{n} \cdot \frac{\qpi_i - 1/n}{\qpi_i + 1/n},
  \end{align*}
  where the last inequality follows because $\card{P} \leq n - 1$.
  And now taking expectation over $\Perm(X)$ and $\Pi$:
  \begin{align*}
    \expect_{\rPerm(X), \rPi}
    \Exp_{\rPerm \given \rPi, \rPerm(X)}{\score(\estar) - \score(e)}
    &=
    \Exp*_{\rPerm(X), \rPi}{
    \Exp_{\rPerm \given \rPi, \rPerm(X)}{\score(\estar)}
    - \expect_{\rPerm(Y) \given \rPi, \rPerm(X)} \Exp_{\rPerm \given \rPerm(X), \rPerm(Y)}{\score(e)}
    }\\
    &\geq
    \Exp*_{\rPerm(X), \rPi}{\sum_{i \in P} \frac{ \paren{\qpi_i - 1/n}^2 }{\qpi_i + 1/n} - 1/n} \\
    &=
    \Exp*_{\rPerm(X), \rPi}{\PTD{\qpi}{\uni_{\Perm(X)}}} - 1/n
    = \rPTD - 1 / n \geq \rPTD / 2.
    \intertext{where the last inequality is because
      $\rPTD \geq \eps \geq 8 / m = 2 / n$.
    }
  \end{align*}
\end{proof}

\begin{claim}\label{clm:variance-score-target}
  The variance of $\score(\estar)$ is at most $\rPTD$.
\end{claim}

\begin{proof}
  Recall that $\Var{X} \leq \Exp{X^2}$ for any random variable $X$.
  As before, we will bound $\Exp{\score(\estar)}$ by conditioning on
  fixed choices of $\Perm(X)$ and the transcript $\Pi$, which
  gives Alice the distribution $\qpi$.
  Also as before, let $P$ denote the set of indices $i$ such that
  $q_i \geq 1 / n$.
  \begin{align*}
    \Exp*_{\rPerm \mid \Perm(X), \Pi}{{\score(\estar)}^2}
    &= \sum_{i \in P} \qpi_i \cdot \pparen*{\frac {\qpi_i - 1 / n}{\qpi_i + 1 / n}}^2\\
    &= \sum_{i \in P} \frac {\qpi_i}{\qpi_i + 1 / n} \cdot \frac{\pparen{\qpi_i - 1/n}^2}{\qpi_i + 1 / n}\\
    &\leq \PTD{\qpi}{ \uni_{\Perm(X)}}.
  \end{align*}
  And taking expectation over $\Perm(X)$ and $\Pi$:
  \[
    \expect_{\rPerm(X), \rPi}
    \Exp*_{\rPerm \mid \rPerm(X), \rPi}{{\score(\estar)}^2}
    \leq \Exp*_{\rPerm(X), \rPi}{\PTD{\qpi}{\uni_{\rPerm(X)}}}
    = \rPTD.
  \]
\end{proof}

\begin{claim}\label{clm:variance-score-other}
  For any $e (\neq \estar) \in X$, the variance of $\score(e)$ is at most
  $2\rPTD$.
\end{claim}

\begin{proof}
  Again, because $\Var{X} \leq \Exp{X^2}$ for any random variable $X$, and
  conditioning on fixed choices of $\Perm(X), \Perm(Y)$ and $\Pi$,
  $\rPerm(e)$ for $e \in X \setminus \Set{\estar}$ is uniformly random over
  $\Perm(X \setminus \Set{\estar})$:
  \begin{align*}
    \Exp*_{\rPerm \mid \Perm(X), \Perm(Y), \Pi}{{\score(e)}^2}
    &= \sum_{i \in P} \frac 1{n - 1}
    \cdot \pparen*{\frac {\qpi_i - 1 / n}{\qpi_i + 1 / n}}^2\\
    &\leq \sum_{i \in P} \frac 1{n - 1}
    \cdot \frac{\paren{\qpi_i - 1/n}^2}{\qpi_i + 1 / n} \cdot \frac 1{1 / n}\\
    &= \PTD*{\qpi}{\uni_{\Perm(X)}} \cdot \pparen*{1 + \frac 1{n - 1}}
    \leq 2\PTD*{\qpi}{ \uni_{\Perm(X)}},
  \end{align*}
  where the first inequality holds because $\qpi_i \geq 0$.
  And taking expectation over $\Perm(X), \Perm(Y)$ and $\Pi$:
  \[
    \Exp*_{\rPerm(X), \rPi}{
      \Exp*_{\rPerm(Y) \mid \rPerm(X), \rPi} {
        \Exp*_{\rPerm \mid \rPerm(X), \rPerm(Y), \rPi}{{\score(e)}^2}
      }
    }
    \leq
    \Exp*_{\rPerm(X), \rPi}{
      2\PTD*{\qpi}{\uni_{\rPerm(X)}}
    }
    = 2\rPTD.
  \]
\end{proof}

\subsubsection{The actual protocol}\label{subsec:actual-protocol}

We are ready to describe the protocol $\pi$.
Let $\gamma \in (0,1)$ be any constant.
We will run $k = 1600 / (\eps \gamma^2)$ (independent) rounds of the
protocol in \Cref{subsec:single-round}, and for each element $e \in X$,
Alice will add up its scores in each round to obtain the total score $\totscore(e)$. \Cref{clm:exp-score-diff} tells us that we should expect the total score for $\estar$ to be higher than the total score of an element $e \neq \estar$.
We then choose the threshold
\[
  \tau = \paren{\Exp{\totscore(\estar)} + \Exp{\totscore(e)}} / 2,
\]
to be exactly in the middle of these two expected values.

Alice will then send to Bob all elements in $X$ which have $\totscore$ at least
$\tau$; let us call this set $S$.
Bob will then compute the intersection $S \cap Y$, and report the unique element
that is (hopefully) present there as the answer.

To complete the reduction, we need to show that:
\begin{itemize}
  \item The total score of the target element $\totscore(\estar)$ is at least
    $\tau$.
    This will ensure that $\estar \in S$, which means that $\pi$ is correct.
  \item The number of elements Alice sends is \emph{small}.
    This is required to control the information cost of $\pi$.
\end{itemize}
We will get both of these properties only with some (constant) high probability;
in particular, if $S$ is too large, Alice will just report failure instead
of communicating $S$.

\begin{claim}
  $\totscore(\estar) \leq \tau$ with probability at most $1 / 100$.
\end{claim}

\begin{proof}
Using Chebyshev's inequality on the total score of $\estar$, we have:

\begin{align*}
  &\mkern-36mu\Prob*{ \card*{\totscore(\estar) - \Exp{\totscore(\estar)}}
  > 1/2 \cdot \Exp{\totscore(\estar) - \totscore(e)} } \\
  &\leq \frac{ \Var*{\totscore(\estar)} }{\pparen*{1/2 \cdot \Exp{\totscore(\estar) - \totscore(e)}}^2}
  = \frac{k \cdot \Var{\score(\estar)}}{k^2 / 4 \cdot \Exp{\score(\estar) - \score(e)}^2}\\
  &\leq \frac{4\rPTD}{k\cdot\pparen{\rPTD/2}^2}
  \leq 16 / k\eps = \gamma^2/100,
\end{align*}
where the second inequality follows from \Cref{clm:exp-score-diff,clm:variance-score-target}.
Hence the probability that $\totscore(\estar)$ is less than $\tau$ is at most
$1/ 100$.
\end{proof}

\begin{claim}
  For any non-$\estar$ element $e \in X$, $\totscore(e) \geq \tau$ with
  probability at most $\gamma^2 / 50$.
\end{claim}

\begin{proof}
Using Chebyshev on the total score of $e \neq \estar$, we have:

\begin{align*}
  &\mkern-36mu\Prob*{ \card*{\totscore(e) - \Exp{\totscore(e)}}
  > 1/2 \cdot \Exp{\totscore(\estar) - \totscore(e)} } \\
  &\leq \frac{ \Var*{\totscore(e)} }{\pparen*{1/2 \cdot \Exp{\totscore(\estar) - \totscore(e)}}^2}
  = \frac{k \cdot \Var{\score(e)}}{k^2 / 4 \cdot \Exp{\score(\estar) - \score(e)}^2}\\
  &\leq \frac{4 \cdot 2\rPTD}{k \cdot \paren{\rPTD / 2}^2}
  \leq 32 / k\eps = \gamma^2/50,
\end{align*}
where the second inequality follows from \Cref{clm:exp-score-diff,clm:variance-score-other}.
\end{proof}

This means that the expected size of $S \setminus \Set{\estar}$ is at most
$\gamma^2 m / 50$, and hence by Markov's this number is $> \gamma^2 m / 10$
with probability $\leq 1/5$.
If the number of elements above the threshold exceeds $\gamma^2 m / 10 + 1$,
then the protocol $\pi$ fails.

And now we have everything we need to bound the information cost of $\pi$.
Note that this part of the argument is exactly as in~\cite{AssadiR20}.
By the chain rule of mutual information, the information cost of $\pi$
on $\dsi$ can be broken down into three components:
\begin{itemize}
  \item The mutual information
    $\Inf{\Perm_1, \ldots , \Perm_k, R_1, \ldots , R_k ; \rX \given \rY}$
    and the symmetric term with $\rX$ and $\rY$ exchanging roles,
    where $\Perm_i$ denotes the random permutations used to scramble the input
    in the $i-$th round of $\pi$, and $R_i$ denotes the randomness of the
    $i$-th run of $\pisi$.
    These are both $0$, because the permutations and the randomness are chosen
    independently of the inputs $\rX$ and $\rY$.
  \item The mutual information
    $\Inf{\rPi_1, \ldots , \rPi_k ; \rX \given \rY, \Sigma, \rR}$
    (and the symmetric term with $\rX$ and $\rY$ swapped), where
    $\rPi_1, \ldots , \rPi_k$ denote the transcripts of the $k$ runs
    of $\pisi$, $\Sigma = \Perm_1, \ldots , \Perm_k$, and
    $\rR = R_1, \ldots , R_k$.
    By the chain rule of mutual information, we have:
    \[ \Inf{\rPi_1, \ldots , \rPi_k ; \rX \given \rY, \Sigma, \rR}
      \leq \sum_{i = 1}^k \Inf{\rPi_i ; \rX \given \rY, \Sigma, \rR, \rPi_{< i}}.
    \]
    Observe that conditioning on $\rX$, $\rY$, $\Perm_i$, and $R_i$ fixes
    $\rPi_i$, and hence makes it independent of $\rPi_{< i}$, $\Perm_{-i}$,
    and $R_{-i}$.
    Applying \Cref{prop:info-decrease}, we get that the $i$-th term of the
    sum above is upper bounded by $\Inf{\rPi_i ; \rX \given \rY, \Perm_i, R_i}$.
    But now (once again) $\rPi_i$ is fixed after conditioning on $\Perm_i(\rX)$,
    $\Perm_i(\rY)$, and $R_i$, and hence independent of the remaining parts
    of $\Perm_i$ and $\rX$ and $\rY$.
    So by \Cref{prop:info-decrease}, each term in the sum is bounded above by
    $\Inf{\rPi_i ; \Perm_i(\rX) \given \Perm(\rY), R_i}$.
    Repeating this entire argument for the symmetric term tells us that:
    \[ \Inf{\rPi_1, \ldots , \rPi_k ; \rX \given \rY, \Sigma, \rR} +
      \Inf{\rPi_1, \ldots , \rPi_k ; \rY \given \rX, \Sigma, \rR}
      \leq \sum_{i = 1}^k
      \Inf{\rPi_i ; \Perm_i(\rX) \given \Perm(\rY), R_i} +
      \Inf{\rPi_i ; \Perm_i(\rY) \given \Perm(\rX), R_i},
    \]
    which is exactly $k \cdot \IntIC_{\dsi}(\pisi)$.
  \item The mutual information between the final set $\rS$ Alice sends and
    her input.
    This is trivially upper bounded by $\Ent{\rS}$, which (since $\rS$ is a
    set of size at most $\gamma^2 \cdot m / 10$) is upper bounded by
    $H_2(\gamma^2 / 10) \cdot m$.
    Since $H_2(p) \leq 2\sqrt{p \cdot (1 - p)} \leq 2\sqrt{p}$
    (by \Cref{fact:ent-ub}),
    this is at most $\gamma m$.
\end{itemize}
So letting $c = \frac{1600}{\gamma^2}$, we get that $\IntIC_{\dsi}(\pi) \le \frac{c}{\eps} \IntIC_{\dsi}(\pisi) + \gamma m$.
\end{proof}

\begin{remark}
  The protocol $\pi$ actually works in the worst case, that is, it needs no
  assumptions about $X$ and $Y$ other than $\card{X} = \card{Y} = m / 4$, and
  that $\card{X \cap Y} = 1$.
  But of course without a distribution on its inputs we cannot talk about
  its information cost.
\end{remark}





\subsection{\texorpdfstring{Short Protocols for $\bhpc$ Give Low-Information $\eps$-Solvers for $\SI$}{Short Protocols for BHCP Give Low-Information eps-Solvers for Set Intersection}}\label{sec:eps-solving-to-bhpc}



We also wish to prove a tight lower bound for $\bhpc_{m,r}$. Note that this lower-bound is not necessary to derive our lower-bounds for $k$-cores and degeneracy, so this section may be skipped for that purpose. The lower-bound ultimately follows from the following lemma, which is analogous to \Cref{lem:mhpc-dist}.

\begin{lemma}\label{lem:nonmhpc-dist}
There is a constant $c > 0$ such that the following statement holds for all $j
\in [r]$:\footnote{Intuitively, the statement says that, if $\pi_\bhpc$
communicates little, then the distribution of $\rZ_j$, conditioned on the
information available at the end of round $j$, is close to uniform.}
\[
\mathsf{Hypothesis}_j:\;\;\;\;\;\;\;\;\E_{(\rE_j, \rPi_j)}\left[\Lambda(\mu(\rZ_j \mid \rE_j, \rPi_j), \mu(\rZ_j))\right] \leq j \cdot c \cdot \left( \frac{CC(\pi_\bhpc)+r\log m}{m^2} + \frac{r} m\right).
\] 

\end{lemma}



One can already see that the statement is more complicated. This is owed to the possibility of the pointers pointing either to some of the (few) instances about which the protocol has revealed information, and the possibility of the pointers looping. These two obstacles raise technical difficulties, which are overcome in two ways: the number of jumps $r$ is assumed to be smaller than $\sqrt m$, so looping is unlikely, and there are only $O(m/r)$ instances about which the protocol could have revealed significant information, hence when hiding an instance of set-intersection in a random position, as we did before, it is unlikely we will hit such an instance.

As before, this is proven by induction. The base case $\mathsf{Hypothesis}_1$ is easy to see: note that $\mu(\rZ_1 \mid \rE_1, \rPi_1) = \mu(\rZ_1)$ as $\rZ_1$ is independent of $\rE_1 = z_0$ which is fixed and $\rZ_1$ is independent of $\rPi_1$ as $\Pi_1$ is a function if $(C,D)$. We obtain
$\mathsf{Hypothesis}_j$ from $\mathsf{Hypothesis}_{j-1}$ in two steps: we first show that if the protocol reveals a lot of information about the variable $\rZ_{j}$, then it must essentially be ``almost-solving'' a \textit{uniformly random instance}. This is because $Z_{j} = T_{Z_{j-1}}$ and $\rZ_{j-1}$ is, under $\mathsf{Hypothesis}_{j-1}$, close to uniform. We then show that this is impossible for a short protocol to do, namely: 

%\begin{lemma}[Progress Lemma] \label{lem:progress}
\begin{restatable}%[Progress Lemma]
{lemma}{nonmprogress} \label{lem:nonmprogress}
There is a constant $c' > 0$ such that for any $j \in [r]$, we have:
\begin{align*}
\E_{(\rE_j, \rPi_j)}\E_{x \sim \cU_\cX}\left[\Lambda(\mu(\rT_x \mid \rE_j, \rPi_j), \mu(\rT_x))\right] \leq  c' \cdot \left( \frac{CC(\pi_\bhpc)+r\log m}{m^2} + \frac{j} m\right)\\
\E_{(\rE_j, \rPi_j)}\E_{y \sim \cU_\cY}\left[\Lambda(\mu(\rT_y \mid \rE_j, \rPi_j), \mu(\rT_y))\right] \leq c' \cdot \left( \frac{CC(\pi_\bhpc)+r\log m}{m^2} + \frac{j} m\right).
\end{align*}
\end{restatable}

We state the proof of \Cref{lem:nonmprogress} shortly. Before that, we show how \Cref{lem:nonmprogress} lets us complete the induction in order to prove \Cref{lem:nonmhpc-dist}. This also follows very closely to the analogous proof of \cite{AssadiCK19} where they showed it for total variation distance.

\begin{proof}[Proof of \Cref{lem:nonmhpc-dist}]
The information---measured by $\Lambda$-distance to uniform---that the $j$-th round of the protocol reveals about $\rZ_j = \rT_{\rZ_{j-1}}$ is the information revealed about an instance sampled according to $\rZ_{j-1}$:
  \begin{align*}
		\E_{(\rE_j,\rProt_j)} \Bracket{\Lambda(\mu(\rZ_j \mid \rE_j,\rProt_j),\mu(\rZ_j))} &=\E_{(\rE_j,\rProt_j)} \Bracket{\Lambda(\mu(\rZ_j \mid \rE_j),\mu(\rZ_j))}\\
		&\hspace{-1.5cm}= \E_{(\rZ^{<j},\rProt^{<j})} \Bracket{\Lambda(\mu(\rZ_j \mid \rZ^{<j},\rProt^{<j}),\mu(\rZ_j))} \tag{by definition of $\rE_j := (\rZ^{<j},\rProt^{<j})$, and since $\rZ_j$ and $\Pi_j$ are independent given $\rE_j$.} \\
		&\hspace{-1.5cm}=\E_{(\rZ^{<j},\rProt^{<j})} \Bracket{\Lambda(\mu(\rT_{\rZ_{j-1}} \mid \rZ^{<j-1},\rZ_{j-1},\rProt^{<j}),\mu(\rZ_j))}\tag{by definition, the pointer $Z_j = T_{Z_{j-1}}$}\\
		&\hspace{-1.5cm}= \E_{(\rZ^{<j-1},\rProt^{<j})} \; \E_{z \sim \rZ_{j-1} \mid \paren{\rZ^{<j-1},\rProt^{<j}}}\; \Bracket{\Lambda(\mu(\rT_{z} \mid \rZ^{<j-1},\rProt^{<j}),\mu(\rZ_j))}.
	\end{align*}
	We removed the conditioning on $\rZ_{j-1}$ in the first distribution because, conditioned on knowing $\rZ^{<j-1},\Prot^{<j}$, for any fixed value of $z$ the value of $\rT_{z}$ is independent of the value of $\rZ_{j-1}$: \begin{itemize}
	    \item If $j-1$ is odd, $\rT_{z}$ is a function of  $(\rC,\rD)$ and $\rZ_{j-1}$ is a function of $(\rA,\rB)$.
            \item If $j-1$ is even, $\rT_{z}$ is a function of $(\rA,\rB)$ and $\rZ_{j-1}$ is a function of $(\bC,\bD)$.
            \item And, furthermore, $(\rA,\rB)$ and $(\rB,\rD)$ are independent given $\rZ^{<j-1}, \rPi^{<j}$. 
	\end{itemize}

    
	\medskip\noindent
    \Cref{lem:nonmprogress} says that if the protocol were to reveal information about the solution $\rT_z$ of a truly uniform instance $z$, then this would give us a protocol for $\eps$-solving set disjointness. But, by induction, we know that the chosen instance $z \sim \rZ_{j-1}$ is close to uniform. It should then follow from \Cref{itm:lambda-loss} that if the protocol reveals some information about this almost uniform instance, it must also reveal information about a truly uniform instance. Let us now make this precise.
    
	Without loss of generality suppose that $j-1$ is odd and hence $z \in \cY$. Then 
	we can upper bound the last expression above by:
	\begin{align*}
		&\E_{(\rE_j,\rProt_j)} \Bracket{\Lambda(\mu(\rZ_j \mid \rE_j,\rProt_j),\mu(\rZ_j))}  \\
  %
		&\qquad \qquad \leq \E_{(\rZ^{<j-1},\rProt^{<j})} \bracket{6 \cdot \E_{\paren{z \sim \cU_{\cY}}}\Bracket{\Lambda(\mu(\rT_{z} \mid \rZ^{<j-1},\rProt^{<j}),\mu(\rZ_j))}} \\
  %
		 &\qquad \qquad  \qquad \qquad \qquad  \qquad + \E_{(\rZ^{<j-1},\rProt^{<j})} \Bracket{\Lambda(\mu(\rZ_{j-1} \mid \rE_{j-1},\rProt_{j-1}),\cU_\cY)} %
   \tag{by \Cref{itm:lambda-loss} and linearity of expectatiom}
   \\
   %
   &\qquad \qquad = 6 \cdot\E_{(\rZ^{<j-1},\rProt^{<j})} \bracket{ \E_{\paren{z \sim \cU_{\cY}}}\Bracket{\Lambda(\mu(\rT_{z} \mid \rZ^{<j-1},\rProt^{<j}),\mu(\rZ_j))}} \tag{$\ast$}\label{eq:final-inequality}\\
		 &\qquad \qquad  \qquad \qquad \qquad  \qquad + \E_{(\rZ^{<j-1},\rProt^{<j})} \Bracket{\Lambda(\mu(\rZ_{j-1} \mid \rE_{j-1},\rProt_{j-1}),\mu(\rZ_{j-1}))}\tag{by definition $z_{j-1} \in \cY$ and $\mu(\rZ_{j-1}) = \cU_\cY$.}\\
   %
   &\qquad \qquad\leq 6 \cdot \underbrace{c' \cdot \left( \frac{CC(\pi_\bhpc)+r\log m}{m^2} + \frac{j} m\right)}_{\text{\Cref{lem:nonmprogress}}} + \underbrace{(j-1) \cdot c \cdot \left( \frac{CC(\pi_\bhpc)+r\log m}{m^2} + \frac{r} m\right)}_{\mathsf{Hypothesis}_{j-1}}\\
		%
  & \qquad \qquad\leq j \cdot c \cdot \left( \frac{CC(\pi_\bhpc)+r\log m}{m^2} + \frac{r} m\right) \\
  &\tag{where we replaced $j \leq r$ by $r$ in the first term and set $c = 6c'$}
	\end{align*}

This concludes the proof.
\end{proof}


\bigskip\noindent
We now prove \Cref{lem:nonmprogress}. The proof shows that any protocol that reveals information about the solution $T_z$ to random instance $z$ can be converted to a low-information protocol for externally $\eps$-solving Set-Intersection. This will later be shown to be impossible. Let us state the lemma again:

\nonmprogress*
Suppose towards a contradiction that this equation does not hold. We use $\protHPC$ to design a protocol $\protSI$
	that can externally $\eps$-solve the Set-Intersection problem $(A_x,B_x)$ for a uniformly at random chosen $x \in \cX$ and appropriately chosen $\eps \in (0,1)$ to be determined later (see \Cref{def:eps-solving} for the notion of $\eps$-solving). 
	
	\begin{tbox}
    \textbf{Protocol $\protSI$}: The protocol for $\eps$-solving $\SI$ using a protocol $\protHPC$ for $\bhpc_{m,r}$. 
    
     \smallskip
    
    \textbf{Input:} An instance $(X, Y) \sim \dsi$ over the universe $\cY$ and
    a number $j \in [r]$.
    
   % \algline
    
    \begin{enumerate}
    	\item \textbf{Sampling the instance.} Alice and Bob create an instance
        $(A, B, C, D)$ of $\bhpc_{m,r}$ as follows: %The players create an instance $(A,B)$ of $\SI$ as follows:
	\begin{enumerate}
	\item Using \underline{public coins}, Alice and Bob sample an index $i \in [m]$ uniformly at random, and  
	Alice sets $A_{x_i} = X$ and Bob sets $B_{x_i} = Y$ using their given inputs in $\SI$. 
	\item Using \underline{public coins}, Alice and Bob sample $A_{x_1}, \ldots, A_{x_{i-1}}$ from the A-side marginal of $\cD_\SI$, and they sample $B_{x_{i+1}}, \ldots, B_{x_{m}}$ from the B-side marginal of $\cD_\SI$. 
	\item Using \underline{private coins}, Alice samples $A_{x_r}$ for $r \in \{i+1, \ldots, n\}$ so that $(A_{x_r},B_{x_r}) \sim \cD_\SI$. Similarly Bob samples $B_{x_\ell}$ for $\ell \in \{1, \ldots, i-1\}$. This completes construction of $(\bA,\bB)$. 
  \item Using \underline{public coins}, Alice and Bob sample $(C, D)$ completely
    from $\cD_\bhpc$ (this is possible by Observation~\ref{obs:distHPC} as
    $(\rA, \rB) \perp (\rC, \rD)$). 
        \end{enumerate}
  
        
           \item\label{line:AB} \textbf{Computing the answer.} 
          For $q = 1$ to $j$, Alice and Bob do the following: \begin{enumerate}
              \item \textbf{Running protocol:} In round $q$, Alice and Bob run
                the round $q$ of the protocol $\prot_\bhpc$ on the instance $(A,
                B, C, D)$ by Alice playing $P_A$, Bob playing $P_B$, and both
                Alice and Bob simulating $P_C$ and $P_D$ with no communication
                (this is possible as both Alice and Bob know $(C, D)$ entirely).
                Denote this transcript by $\Pi_q$.

           \item \textbf{Fixing pointer:} Alice and Bob compute the pointer $z_{q}$ using the fact
           that for any underlying instance $(A_x,B_x) \in (A, B) \setminus (A_{x_i},B_{x_i})$ either Alice or Bob knows the entire instance. We use $\Prot^*$ to denote the cumulative transcript generated by this step in all rounds.
           
           \item They stop if $x_i = z_{q}$.
          \end{enumerate} 
           
           
	\item The players return $\Pi_\SI := (\Pi_1,\ldots,\Pi_j,\Pi^*)$ from which $\mu(\estar)$ is calculated.
    \end{enumerate}
\end{tbox}

First, we show that the internal information cost of this protocol is small if the communication cost of $\Pi_\bhpc$ is small. This uses standard direct sum argument of information complexity. More concretely, we have the following:

\begin{claim} \label{clm:nonmdirect-sum}
    $\ICost{\pi_\SI}{\cD} = O\left(\frac{CC(\pi_\bhpc)}{m} + \frac{j \log m} m\right)$.
\end{claim}



The proof of \Cref{clm:nonmdirect-sum} is similar to that of Claim 5.6 of~\cite{AssadiCK19}. We defer it to the \Cref{sec:proof-direct-sum}. At this point, we also need a bound on how well $\Pi_\SI$ can ``almost'' solve Set-Intersection. To this end, we have the following guarantee.

\begin{claim} \label{clm:nonmaverage-coordinate}
$\pi_\SI$ externally $\eps$-solves Set-Intersection on $\cD$ where:
\[
  \eps \geq \E_{(\rE_j, \rPi_j)}\E_{x \sim \cU_\cX}
  \left[\Lambda(\mu(\rT_x \mid \rE_j, \rPi_j), \mu(\rT_x)) \right] - \frac{j} m.
\]
\end{claim}



\begin{proof} 
    This proof follows the proof of Claim 5.5 of \cite{AssadiCK19} almost verbatim. The only change is that we are working with positive triangular discrimination ($\Lambda$) instead of the total variation distance ($\Delta_{\sf TVD}$). However, as $\Lambda$ is bounded by $\Delta_{\sf TVD}$ (\Cref{prop:prop-lambda}), a similar calculation yields the claim.
    By \Cref{def:eps-solving}, $\eps = \E_{\rPi_\SI}[\Lambda(\mu(\rT \mid
    \rPi_\SI), \mu(\rT))]$. So we have:

    \begin{align*}
        \E_{\rPi_\SI}[\Lambda(\mu(\rT \mid \rPi_\SI), \mu(\rT))]
        &= \E_{\rE_j, \rPi_j, \rPi^\ast, \rI}
        [\Lambda(\mu(\rT_{x_{\rI}} \mid \rPi_\SI), \mu(\rT_{x_{\rI}}))]
        \tag{as $\rT = \rT_{x_{\rI}}$}\\
        &\hspace{-1.9cm}=
        \E_{\rE_j, \rPi_j} \E_{\rI} \E_{\rPi^\ast\mid \rE_j, \rPi_j, \rI}
        [\Lambda(\mu(\rT_{x_{\rI}} \mid \rPi_\SI), \mu(\rT_{x_i}))]
        \tag{as $\rI$ is independent of $(\rE_j, \rPi_j)$}\\
        &\hspace{-1.9cm}=
        \E_{\rE_j, \rPi_j}
        \left[\sum_i \frac 1 m \E_{\rPi^\ast\mid \rE_j, \rPi_j, \rI = i}
        [\Lambda(\mu(\rT_{x_i} \mid \rPi_\SI), \mu(\rT_{x_i}))]\right]
        \tag{as $\rI \sim \cU_{[m]}$}\\
        &\hspace{-1.9cm}=
        \E_{\rE_j, \rPi_j} \left[\sum_{x_i \in \rZ^{<j}}
          \frac 1 m \Lambda(\mu(\rT_{x_i} \mid \rPi_\SI), \mu(\rT_{x_i}))
        + \sum_{x_i \notin \rZ^{<j}} \frac 1 m \Lambda(\mu(\rT_{x_i} \mid \rE_j,
      \rPi_j), \mu(\rT_{x_i}))\right]\\
        \tag{If $x_i \notin \rZ^{< j}$ then $\rPi_\SI = (\rE_j, \rPi_j)$}\\
         &\hspace{-1.9cm}\geq
         \E_{\rE_j, \rPi_j} \left[\sum_{x_i \notin \rZ^{<j}} \frac 1 m
         \Lambda(\mu(\rT_{x_i} \mid \rE_j, \rPi_j), \mu(\rT_{x_i}))\right]
         \tag{as the first sum term is at least 0}\\
         &\hspace{-1.9cm}\geq \E_{\rE_j, \rPi_j} \E_i\left[
         \Lambda(\mu(\rT_{x_i} \mid \rE_j, \rPi_j), \mu(\rT_{x_i}))\right]
         -\frac{j}{m}. \tag{as $\Lambda$ is upper bounded by 1}
    \end{align*}
\end{proof}



\noindent
Using \Cref{clm:nonmdirect-sum} and the lower bound on $\ICost{\pi_\SI}{\cD}$ (\Cref{cor:si-ic-lb-external}), we have:
\[
\eps = O\left(\frac{CC(\pi_\bhpc)}{m^2} + \frac{j \log m} {m^2}\right)
\]
Now, using \Cref{clm:nonmaverage-coordinate}, we get:
\[
\E_{(\rE_j, \rPi_j)}\E_{x \sim \cU_\cX} \left[\Lambda(\mu(\rT_x \mid \rE_j,
\rPi_j), \mu(\rT_x)) \right] - \frac j n = O\left(\frac{CC(\pi_\bhpc)}{m^2} + \frac{j \log m} {m^2}\right),
\]
which concludes the proof of \Cref{lem:nonmprogress}.


\subsection{Putting it All Together}\label{sec:put-together}


 Given \Cref{lem:mhpc-dist}, we prove \Cref{thm:lb-mhpc} in the following way:
 Consider a protocol $\pi$ that is solving $\bmhpc_{m,r}$ with communication
 cost $o(m^2)$.
 By \Cref{lem:mhpc-dist}, we have:

 \begin{align*}
     \E_{(\rE_r, \rPi_r)}\left[\Lambda(\mu(\rZ_r \mid \rE_r, \rPi_r), \mu(\rZ_r))\right]
     &\leq 6 c \cdot \frac{\mi{\rPi}{\rA,\rC \mid \rB,\rD} + \mi{\rPi}{\rB,\rD
\mid \rA,\rC}}{m^2}\\ &\underset{(\ast)}{\leq} 12c \cdot \frac{\myCC(\pi_\bhpc)}{m^2} = o(1),
 \end{align*}

where inequality $(\ast)$ follows from \Cref{part:uniform,part:info-atmost-rv}
in \Cref{fact:it-facts}.
However, next, we show that, for a correct protocol for $\bmhpc_{m,r}$, this
quantity should be $\Omega(1)$, which contradicts the assumption that
$\myCC(\pi_\bhpc) = o(m^2)$.
An entirely similar argument derives \Cref{thm:lb-hpc} from \Cref{lem:nonmhpc-dist}.
The only difference between these two arguments is the bound on $r$: We have to
assume $r = o(\sqrt{m})$ because of the additive $r^2/m$ factor in
\Cref{lem:nonmhpc-dist} (when applied with on the $r$-th round of the protocol).




\begin{lemma} \label{lem:hpc-correctness}
    Any protocol that solves $\bmhpc_{m,r}$ on distribution $\cD_\bmhpc$, or $\bhpc_{m,r}$, on distribution $\cD_\bhpc$, with constant probability must have:
    \[
 \E_{(\rE_r, \rPi_r)}\left[\Lambda(\mu(\rZ_r \mid \rE_r, \rPi_r), \mu(\rZ_r))\right] = \Omega(1).
\]
\end{lemma}

\begin{proof}
    
 Let $b$ denote the modulus function such that $b(z_k) = i \mod 2$ iff $z_r = x_i$ or $y_i$ depending on where $r$ is even or odd. If the protocol $\pi$ computes $\bmhpc_{m,r}$ or $\bhpc_{m,r}$ with probability $2/3$, we have the following inequality:

 \begin{align*}
     2/3 &= \E_{\rE_r,\rPi_r} \Pr_{\rZ_r \mid \rE_r,\rPi_r} [b(\rZ_r) =
     \cO_\pi(\rE_r,\rPi_r)]\\
     &\underset{(1)}\leq \E_{\rE_r,\rPi_r}\left[\Pr_{\rZ_r \sim
     \cU_{\cX}}[b(\rZ_r) = \cO_\pi(\rE_r,\rPi_r)] + \|\mu(b(\rZ_r) \mid \rE_r,\rPi_r) - \cU_{0,1}\|_1 \right]\\
     &\underset{(2)}\leq \E_{\rE_r,\rPi_r}\left[\Pr_{\rZ_r \sim
     \cU_{\cX}}[b(\rZ_r) = \cO_\pi(\rE_r,\rPi_r)] + \|\mu(\rZ_r \mid \rE_r,\rPi_r) - \cU_\cX\|_1 \right]\\
     &= \frac 1 2 + \E_{\rE_r,\rPi_r}\left[\|\mu(\rZ_r \mid \rE_r,\Pi_r) - \cU_\cX\|_1 \right]
 \end{align*}
 where (1) holds because of the simple fact that $\Pr_p[\cE] \leq \Pr_q[\cE] + \|p-q\|_1$ where $p$ and $q$ are two distributions over a sample space of which $\cE$ is an event (see Footnote \ref{fn:tvd}), and (2) holds because of data processing inequality. This implies:
 \begin{align*}
     \E_{\rE_r,\rPi_r}\left[\|\mu(\rZ_r \mid \rE_r,\rPi_r) - \cU_{0,1}\|_1
     \right] \geq 1/6 &\Rightarrow \left(\E_{\rE_r,\rPi_r}\left[\|\mu(\rZ_r \mid
     \rE_r,\rPi_r) - \cU_{0,1}\|_1 \right] \right)^2\geq 1/36\\
     &\Rightarrow \E_{\rE_r,\Pi_r}\left[\|\mu(\rZ_r \mid \rE_r,\rPi_r) -
     \cU_{0,1}\|_1^2 \right] \geq 1/36\\
     &\Rightarrow \E_{\rE_r,\Pi_r}
     \left[\Lambda(\mu(\rZ_r \mid \rE_r,\Pi_r), \cU_{0,1}) \right]
     \geq 1/288 = \Omega(1),
 \end{align*}
 which completes the proof.
\end{proof}




\appendix

\section{Sampling Bounds}\label{sec:sampling}

Recall the Chernoff bound:


\begin{proposition}\label{lem:chernoff}
    Let $X_1, \cdots, X_n$ be independent random variables with values in $[0,1]$ and $X = \sum_i X_i$. Then: \begin{itemize}
        \item \textbf{Additive form: } For any $b \ge 0$, $\Pr[|X - \E[X]| \ge b] \le 2\cdot \exp\left(-\frac{b^2}{2n}\right)$, \label{itm:chernoff-add}

        \item \textbf{Multiplicative form: } For $0 \le \delta \le 1$, $\Pr[|X - \E[X]| \ge (1+\delta) E[X]] \le \exp\left(-\frac{\delta^2 E[X]}{3}\right)$.\label{itm:chernoff-mult}
    \end{itemize}
\end{proposition}

\bigskip\noindent
Now suppose we are allowed to repeatedly sample from an unknown Bernoulli distribution $\mu$. We know that $\mu$ either has acceptance probability $p$ or acceptance probability $q$. We wish to determine which is the case, via some method that may err with probability $\gamma$. Suppose that, for some $1 \ge \alpha,\beta > 0$, $p \le \alpha$, whereas $q \ge (1 + \beta) \alpha$. We may then draw 
\[
k = k(\alpha, \beta, \gamma) = \left\lceil\frac{48}{\alpha \beta^2} \ln\left(\frac{1}{\gamma}\right)\right\rceil
\]
samples $X_1, \ldots, X_k$ from $\mu$, and choose the threshold
\[
\tau = \tau(\alpha, \beta, \gamma) = \alpha \cdot \left(1 + \frac{\beta}{4}\right) \cdot k 
\]
It now follows, from the multiplicative form of the Chernoff bound (\Cref{lem:chernoff}), that
\[
\Pr[\sum_i X_i \ge \tau] = \Pr[\sum_i X_i \ge \alpha k \;\left(1 + \frac\beta4\right)] \le \exp\left(- \frac{\beta^2}{48} \alpha k \right) \le \gamma.
\]
Now, whenever $\beta \in [0,1]$, we have $1 + \frac\beta4 \le (1 + \beta) ( 1 - \frac{\beta}{4})$, and so, if the acceptance probability is $q \ge \alpha(1 + \beta)$, then
\begin{align*}
\Pr[\sum_i X_i \le \tau] %
& = \Pr[\sum_i X_i \le \alpha k\left(1 + \frac\beta4\right)] \le \Pr[\sum_i X_i \le \alpha (1 + \beta) k\left(1 - \frac\beta4\right)] \\
& \le \exp\left(- \frac{\beta^2}{48} \alpha (1+\beta) k \right) \le \gamma.
\end{align*}





\section{Proof of Propositions Regarding Triangular Discrimination} \label{sec:td-props}

\lambdatvd*

\begin{proof}
To see the lower bound on $\Lambda(\mu, \nu)$:
\[
\frac{\|\mu - \nu\|_1} 2 = \sum_{x: \mu(x) > \nu(x)} (\mu(x) - \nu(x)) =  \sum_{x: \mu(x) > \nu(x)} \frac{(\mu(x) - \nu(x))^2}{\sqrt{\mu(x) + \nu(x)}} \sqrt{\mu(x) + \nu(x)} \underset{(*)}{\leq} \sqrt{2\Lambda(\mu, \nu)}.
\]
Here $(*)$ follows from Cauchy-Schwarz. To see the upper bound:
\[
\Lambda(\mu, \nu) = \sum_{x: \mu(x) > \nu(x)}\frac{(\mu(x) - \nu(x))^2}{\mu(x) + \nu(x)} = \sum_{x: \mu(x) > \nu(x)}\frac{(\mu(x) - \nu(x))^2}{\mu(x) - \nu(x)} \cdot \frac{\mu(x) - \nu(x)}{\mu(x) + \nu(x)} \underset{(**)}{\leq} \frac{\|\mu - \nu\|_1} 2.
\]
Here $(**)$ follows from $\frac{\mu(x) - \nu(x)}{\mu(x) + \nu(x)}\leq 1$.

\end{proof}

We will make use of the following:

\begin{proposition}[\cite{Yehudayoff20}]
\label{fact:trick1}
If $|\eta| \leq \sqrt{a(b+ \eta)}$ with $a,b\geq 0$, then $\eta \leq a + 2b$.
\end{proposition}

\begin{proof}
    If $a = 0$, then $\eta = 0$ and the upper-bound follows. Otherwise, if $\eta^2 - a \eta - a b \le 0$, then
    \[
    \eta \le \frac{a + \sqrt{a^2 + 4 a b}}{2} = \frac{a}{2} (1 + \sqrt{1 + 4 b/a}) \le \frac{a}{2}(1 + 1 + 4 b/a). \qedhere
    \]
\end{proof}

\lambdaloss*

\begin{proof}
We have:
\begin{align*}
    |\eta|=\left|\E_\mu(f(x)) - \E_\nu(f(x))\right| &= \left|\sum_x(\mu(x) - \nu(x))f(x)\right|\\
    &\leq \sum_x\left|(\mu(x) - \nu(x))\right|f(x)\tag{By triangle inequality}\\
    &= \sum_x\frac{|\mu(x) - \nu(x)|\sqrt{f(x)}}{\sqrt{\mu(x) + \nu(x)}} \sqrt{(\mu(x) + \nu(x)) f(x)}\\
    &\leq \sqrt{\sum_x\frac{(\mu(x) - \nu(x))^2f(x)}{\mu(x) + \nu(x)}} \sqrt{\sum_x(\mu(x) + \nu(x)) f(x)}\tag{By Cauchy-Schwarz}
\end{align*}

\noindent
Let us look at the first term without the square root:
\begin{align*}
    \sum_x\frac{(\mu(x) - \nu(x))^2f(x)}{\mu(x) + \nu(x)} &= \sum_{x: \mu(x) \geq \nu(x)}\frac{(\mu(x) - \nu(x))^2f(x)}{\mu(x) + \nu(x)} + \sum_{x: \mu(x) < \nu(x)}\frac{(\mu(x) - \nu(x))^2f(x)}{\mu(x) + \nu(x)}\\
    &\leq \Lambda(\mu,\nu)\cdot \max_x f(x) + \sum_{x: \mu(x) < \nu(x)}\frac{(\mu(x) - \nu(x))^2f(x)}{\mu(x) + \nu(x)}\\
    &\leq \Lambda(\mu,\nu)\cdot \max_x f(x) + \sum_{x}\frac{(\nu(x))^2}{\nu(x)}f(x) = \Lambda(\mu,\nu)\cdot \max_x f(x) + \E_\nu(f(x)).
\end{align*}
which implies:
\begin{align*}
    |\eta|=\left|\E_\mu(f(x)) - \E_\nu(f(x))\right| \leq \sqrt{\underbrace{(\Lambda(\mu,\nu)\cdot \max_x f(x) + \E_\nu(f(x)))}_{a} (\eta+ \underbrace{2\E_\nu(f(x)))}_b)}
\end{align*}
At this point, we can use \Cref{fact:trick1} to get the following:
\begin{align*}
    \E_\mu(f(x)) - \E_\nu(f(x)) \leq \Lambda(\mu,\nu)\cdot \max_x f(x) + 5\Pr_\nu(\cE),\\
    \text{i.e., }\E_\mu(f(x))\leq \Lambda(\mu,\nu)\cdot \max_x f(x) + 6\E_\nu(f(x)),
\end{align*}
which completes the proof.
\end{proof}

\convex*

\begin{proof}
To see this, it suffices to compute the Hessian of each term.
We have
\[
\Lambda(\mu,\nu) = \sum_{x} \begin{cases}
\frac{(\mu(x)-\nu(x))^2}{\mu(x)+\nu(x)} & \text{if } \mu(x) \ge \nu(x)\\
0 & \text{otherwise}.
\end{cases}
\]
So let us prove that each term is jointly convex in the values of $(\mu(x), \nu(x))$. We restrict our attention to the unit square. The function is zero, hence convex, inside the region where $\mu(x) \le \nu(x)$. Furthermore, it is (obviously) non-negative inside the region where $\mu(x) > \nu(x)$. It is also convex in this region, since the Hessian is positive semidefinite there:
 \[
 \nabla^2 \frac{(\mu(x) - \nu(x))^2}{\mu(x)+\nu(x)}  =
\begin{pmatrix}
 \frac{8 \nu(x)^2}{(\mu(x) +\nu(x) )^3} & -\frac{8 \mu(x)  \nu(x) }{(\mu(x) +\nu(x) )^3} \\
 -\frac{8 \mu(x)  \nu(x) }{(\mu(x) +\nu(x) )^3} & \frac{8 \mu(x)^2}{(\mu(x) +\nu(x) )^3}
\end{pmatrix}.
\]
(The matrix is PSD since the off-diagonal entry, in absolute value, does not exceed the geometric mean of the diagonal entries.) It follows that each term is convex, and hence $\Lambda(\mu,\nu)$ is convex, in $(\mu,\nu)$.
\end{proof}



\section{Proof of \texorpdfstring{\Cref{clm:nonmdirect-sum}}{Claim}} \label{sec:proof-direct-sum}


\begin{proof}
This proof is a verbatim reproduction of that in \cite{AssadiCK19}. For any $i \in [m]$, define $\bA^{<i} := \{A_{x_1},\ldots,A_{x_{i-1}}\}$, $\bB^{>i} := \{B_{x_{i+1}},\ldots,B_{x_{m}}\}$. Recall that 
the internal information cost of $\prot_\SI$ is $\IC{\prot_\SI}{\cD_\SI} := \mi{\rA}{\rProtSI \mid \rB} + \mi{\rB}{\rProtSI \mid \rA}$. In the following, we focus on bounding the first term. The second term can be bounded exactly
the same by symmetry.  

As $(\rI,\bA^{<\rI},\bB^{>\rI},\bC,\bD)$ is sampled via public randomness in $\prot_\SI$, by Proposition~\ref{prop:public-random},   
\begin{align*}
	\mi{\rA}{\rProtSI \mid \rB} &= \mi{\rA}{\rProtSI \mid \rB, \rI,\bA^{<\rI},\bB^{>\rI},\bC,\bD} \leq \mi{\rA}{\rProtSI \mid \rB, \rI,\bA^{<\rI},\bB^{>\rI}}. 
\end{align*}
The inequality is by Proposition~\ref{prop:info-decrease} as we now show $\rA \perp (\bC,\bD) \mid \rProtSI,\rB,\rI,\bA^{<\rI},\bB^{>\rI}$ (and hence conditioning on $(\bC,\bD)$ can only decrease the mutual information). 
This is because $\rA \perp (\bC,\bD) \mid \rB,\rI,\bA^{<\rI},\bB^{>\rI}$ by Observation~\ref{obs:distHPC} and $\rProtSI$ is transcript of a deterministic protocol plus $z_1,\ldots,z_j$ obtained deterministically and hence we can apply
Proposition~\ref{prop:hpc-rectangle}.  

Define a random variable $\rTheta \in \{0,1\}$ where $\rTheta = 1$ iff in Line~(\ref{line:AB}) of protocol $\prot_\SI$, we terminate the protocol. In other words $\rTheta = 1$ iff $x_i \in Z^{<j}$. Since $\rA \perp \rTheta \mid \rB, \rI,\bA^{<\rI},\bB^{>\rI}$, 
further conditioning on $\rTheta$ can only increase the mutual information term above by Proposition~\ref{prop:info-increase}, hence, 
\begin{align}
	\mi{\rA}{\rProtSI \mid \rB} &\leq \mi{\rA}{\rProtSI \mid \rB, \rI,\bA^{<\rI},\bB^{>\rI},\rTheta} \notag \\
	&= \frac{m-j}{m} \cdot \mi{\rA}{\rProtSI \mid \rB, \rI,\bA^{<\rI},\bB^{>\rI},\rTheta=0} + \frac{j}{m} \cdot \mi{\rA}{\rProtSI \mid \rB, \rI,\bA^{<\rI},\bB^{>\rI},\rTheta=1} \notag \\
	&\leq \frac{m-j}{m} \cdot \mi{\rA}{\rProtSI \mid \rB, \rI,\bA^{<\rI},\bB^{>\rI},\rTheta=0}, \label{eq:mi-ra}
\end{align}
since conditioned on $\rTheta=1$, the protocol $\rProtSI$ is simple some prefix of $\rZ^{<j}$ and is hence independent of the input $(\rA,\rB)$ and carries no information about $\rA$ (see \itfacts{info-zero}). 
We now further bound the RHS of Eq~(\ref{eq:mi-ra}). When $\rTheta = 0$, $\rProtSI = (\rZ^{<j} ,\rProt_1,\ldots,\rProt_j) = (\rE^{<j},\rProt_j)$. Hence, we can write, 
  \begin{align*}
	 \mi{\rA}{\rProtSI \mid \rB, \rI,\bA^{<\rI},\bB^{>\rI},\rTheta=0} &\leq \mi{\rA}{\rE_j,\rProt_j \mid \rB, \rI,\bA^{<\rI},\bB^{>\rI},\rTheta=0} \\
	&= \mi{\rA}{\rZ^{<j} \mid \rB, \rI,\bA^{<\rI},\bB^{>\rI},\rTheta=0} \\
	&\hspace{25pt} + \mi{\rA}{\rProt^{<j},\rProt_j \mid \rZ^{<j},\rB, \rI,\bA^{<\rI},\bB^{>\rI},\rTheta=0} \tag{by chain rule in \itfacts{chain-rule} and since $\rE_j = (\rProt^{<j},\rZ^{<j})$} \\
	&\leq \mi{\rA}{\rProt \mid \rZ^{<j},\rB, \rI,\bA^{<\rI},\bB^{>\rI},\rTheta=0},
\end{align*}
as $\rA \perp \rZ^{<j} \mid \rTheta=0$ (and other variables) and hence the first term is zero, and in the second term $\rProt$ contains $\rProt^{<j},\rProt_j$ (plus potentially other terms) and so having $\rProt$ in instead can only increase
the information. By further expanding the conditional information term above,
\begin{align*}
	 &\mi{\rA}{\rProtSI \mid \rB, \rI,\bA^{<\rI},\bB^{>\rI},\rTheta=0} \\
	 &\hspace{50pt}\leq \E_{(Z^{<j},i) \mid \rTheta=0} \bracket{\mi{\rA}{\rProt \mid \rB, \bA^{<i},\bB^{>i},\rI=i,\rZ^{<j}=Z^{<j},\rTheta=0}} \\
	&\hspace{50pt}= \E_{Z^{<j} \mid \rTheta=0} \bracket{\sum_{\substack{i=1 \\~i \notin Z^{<j}}}^{m} \frac{1}{m-j}  \mi{\rA_{x_i}}{\rProt \mid \rB_{x_i}, \bA^{<i},\bB^{>i},\rI=i,\rZ^{<j}=Z^{<j},\rTheta=0}} \tag{conditioned on $\rTheta=0$, $i$ is chosen
	uniformly at random from $Z^{<j}$; also $(\rA,\rB) = (\rA_{x_i},\rB_{x_i})$} \\
	&\hspace{50pt}= \E_{Z^{<j} \mid \rTheta=0} \bracket{\sum_{i \notin Z^{<j}} \frac{1}{m-j} \cdot \mi{\rA_{x_i}}{\rProt \mid \rB_{x_i}, \bA^{<i},\bB^{>i},\rZ^{<j}=Z^{<j},\rTheta=0}}
	\tag{we dropped the conditioning on $\rI=i$ as all remaining variables are independent of this event}  \\
	&\hspace{50pt}= \E_{Z^{<j} \mid \rTheta=0} \bracket{\sum_{i \notin Z^{<j}} \frac{1}{m-j} \cdot \mi{\rA_{x_i}}{\rProt \mid \bA^{<i},\bB,\rZ^{<j}=Z^{<j},\rTheta=0}}
	\tag{as $\rA_{x_i} \perp \bB^{<i} \mid \rB_{x_i},\bA^{<i}$ by Observation~\ref{obs:distHPC} and hence we can apply Proposition~\ref{prop:info-increase}} \\
	&\hspace{50pt}\leq \E_{Z^{<j} \mid \rTheta=0} \bracket{\sum_{i=1}^{m} \frac{1}{m-j} \cdot \mi{\rA_{x_i}}{\rProt \mid \bA^{<i},\bB,\rZ^{<j}=Z^{<j},\rTheta=0}} \\ 
	\tag{mutual information is non-negative by \itfacts{info-zero} and so we can add the terms in $Z^{<j}$  as well} \\
	&\hspace{50pt}= \E_{Z^{<j} \mid \rTheta=0} \bracket{\sum_{i=1}^{m} \frac{1}{m-j} \cdot \mi{\rA_{x_i}}{\rProt \mid \bA^{<i},\bB,\rZ^{<j}=Z^{<j},\rTheta=0}} \\
	&\hspace{50pt}= \frac{1}{m-j} \cdot \E_{Z^{<j} \mid \rTheta=0} \bracket{\mi{\bA}{\rProt \mid \bB,\rZ^{<j}=Z^{<j},\rTheta=0}} \tag{by chain rule in \itfacts{chain-rule}} \\
	&\hspace{50pt}= \frac{1}{m-j} \cdot \mi{\bA}{\rProt \mid \bB,\rZ^{<j},\rTheta=0} \\
	&\hspace{50pt}\leq \frac{1}{m-j} \cdot \paren{\mi{\bA}{\rProt \mid \bB,\rTheta=0} + \en{\rZ^{<j}}} \\
	&\hspace{50pt}= \frac{1}{m-j} \cdot \paren{\mi{\bA}{\rProt \mid \bB} + \en{\rZ^{<j}}} \tag{transcript of the protocol $\prot_\bhpc$ (namely $\rProt$) on input $(\bA,\bB)$ is independent of $\rTheta$} \\
	&\hspace{50pt}\leq \frac{1}{m-j} \cdot \paren{\en{\rProt} + \en{\rZ^{<j}}} \leq \frac{\CC{\prot_\bhpc}{} }{m-j}   + \frac{j \cdot \log{m}}{m-j}. \tag{by sub-additivity of entropy (\itfacts{sub-additivity}) and \itfacts{uniform}}
\end{align*}

By plugging in this bound in Eq~(\ref{eq:mi-ra}), we have that,
\begin{align*}
	\mi{\rA}{\rProtSI \mid \rB} &\leq \frac{m-j}{m} \cdot \paren{\frac{\CC{\prot_\bhpc}{} }{m-j}   + \frac{j \cdot \log{m}}{m-j}} = \frac{\CC{\prot_\bhpc}{} }{m}   + \frac{j \cdot \log{m}}{m}. 
\end{align*}
By symmetry, we can also prove the same bound on $\mi{\rB}{\rProtSI \mid \rA}$. As such, we have,
\begin{align*}
	\mi{\rA}{\rProtSI \mid \rB} + \mi{\rB}{\rProtSI \mid \rA} \leq 2 \cdot \paren{\frac{\CC{\prot_\bhpc}{} }{m}   + \frac{j \cdot \log{m}}{m}}.
\end{align*}

We shall note that strictly speaking the factor $2$ above is not needed (similar to the proof of Proposition~\ref{prop:cc-ic}) but as this factor is anyway suppressed through O-notation later in the proof, the above bound
suffices for our purpose. %\qed\; %\Cref{clm:direct-sum}

\end{proof}

%\clearpage
\section{Background on Information Theory}\label{prelim-sec:info}

We now briefly introduce some definitions and facts from information theory that are needed in this thesis. We refer the interested reader to the text by Cover and Thomas~\cite{CoverT06} for an excellent introduction to this field, 
and the proofs of the statements used in this Appendix. 

For a random variable $\rA$, we use $\supp{\rA}$ to denote the support of $\rA$ and $\distribution{\rA}$ to denote its distribution. 
When it is clear from the context, we may abuse the notation and use $\rA$ directly instead of $\distribution{\rA}$, for example, write 
$A \sim \rA$ to mean $A \sim \distribution{\rA}$, i.e., $A$ is sampled from the distribution of random variable $\rA$. 


We denote the \emph{Shannon Entropy} of a random variable $\rA$ by
$\en{\rA}$, which is defined as: 
\begin{align}
	\en{\rA} := \sum_{A \in \supp{\rA}} \Pr\paren{\rA = A} \cdot \log{\paren{1/\Pr\paren{\rA = A}}} \label{eq:entropy}
\end{align} 

\noindent
The \emph{conditional entropy} of $\rA$ conditioned on $\rB$ is denoted by $\en{\rA \mid \rB}$ and defined as:
\begin{align}
\en{\rA \mid \rB} := \E_{B \sim \rB} \bracket{\en{\rA \mid \rB = B}}, \label{eq:cond-entropy}
\end{align}
where 
$\en{\rA \mid \rB = B}$ is defined in a standard way by using the distribution of $\rA$ conditioned on the event $\rB = B$ in Eq~(\ref{eq:entropy}). 

\noindent
Recall that the binary entropy $H(p)$ is the entropy of a two-valued random
variable that takes one value with probability $p$.
We need the following upper bound on the binary entropy.

\begin{fact}\label{fact:ent-ub}
  The binary entropy $H_2(p) \leq 2 \sqrt{ p \cdot (1 - p) }$.
\end{fact}


 
The \emph{mutual information} of two random variables $\rA$ and $\rB$ is denoted by
$\mi{\rA}{\rB}$ and is defined as:
\begin{align}
\mi{\rA}{\rB} := \en{A} - \en{A \mid  B} = \en{B} - \en{B \mid  A}. \label{eq:mi}
\end{align}
\noindent
The \emph{conditional mutual information} $\mi{\rA}{\rB \mid \rC}$ is $\en{\rA \mid \rC} - \en{\rA \mid \rB,\rC}$ and hence by linearity of expectation:
\begin{align}
	\mi{\rA}{\rB \mid \rC} = \E_{C \sim \rC} \bracket{\mi{\rA}{\rB \mid \rC = C}}. \label{eq:cond-mi}
\end{align} 



\subsection{Useful Properties of Entropy and Mutual Information}\label{sec:prop-en-mi}

We shall use the following basic properties of entropy and mutual information throughout. 


\begin{fact}\label{fact:it-facts}
  Let $\rA$, $\rB$, $\rC$, and $\rD$ be four (possibly correlated) random variables.
   \begin{enumerate}
  \item \label{part:uniform} $0 \leq \en{\rA} \leq \log|\supp{\rA}|$. The right equality holds
    iff $\distribution{\rA}$ is uniform.
  \item \label{part:info-zero} $\mi{\rA}{\rB \mid \rC} \geq 0$. The equality holds iff $\rA$ and
    $\rB$ are \emph{independent} conditioned on $\rC$.
    \item \label{part:info-atmost-rv} $\mi{\rA}{\rB \mid \rC} \leq \en{\rA}$.
  \item \label{part:cond-reduce} \emph{Conditioning on a random variable reduces entropy}:
    $\en{\rA \mid \rB,\rC} \leq \en{\rA \mid  \rB}$.  The equality holds iff $\rA \perp \rC \mid \rB$.
    \item \label{part:sub-additivity} \emph{Subadditivity of entropy}: $\en{\rA,\rB \mid \rC}
    \leq \en{\rA \mid C} + \en{\rB \mid  \rC}$.
   \item \label{part:ent-chain-rule} \emph{Chain rule for entropy}: $\en{\rA,\rB \mid \rC} = \en{\rA \mid \rC} + \en{\rB \mid \rC,\rA}$.
  \item \label{part:chain-rule} \emph{Chain rule for mutual information}: $\mi{\rA,\rB}{\rC \mid \rD} = \mi{\rA}{\rC \mid \rD} + \mi{\rB}{\rC \mid  \rA,\rD}$.
  \item \label{part:data-processing} \emph{Data processing inequality}: for a function $f(\rA)$ of $\rA$, $\mi{f(\rA)}{\rB \mid \rC} \leq \mi{\rA}{\rB \mid \rC}$. 
   \end{enumerate}
\end{fact}


\noindent
We also use the following two standard propositions, regarding the effect of conditioning on mutual information.

\begin{proposition}\label{prop:info-increase}
  For random variables $\rA, \rB, \rC, \rD$, if $\rA \perp \rD \mid \rC$, then, 
  \[\mi{\rA}{\rB \mid \rC} \leq \mi{\rA}{\rB \mid  \rC,  \rD}.\]
\end{proposition}
 \begin{proof}
  Since $\rA$ and $\rD$ are independent conditioned on $\rC$, by
  \itfacts{cond-reduce}, $\HH(\rA \mid  \rC) = \HH(\rA \mid \rC, \rD)$ and $\HH(\rA \mid  \rC, \rB) \ge \HH(\rA \mid  \rC, \rB, \rD)$.  We have,
	 \begin{align*}
	  \mi{\rA}{\rB \mid  \rC} &= \HH(\rA \mid \rC) - \HH(\rA \mid \rC, \rB) = \HH(\rA \mid  \rC, \rD) - \HH(\rA \mid \rC, \rB) \\
	  &\leq \HH(\rA \mid \rC, \rD) - \HH(\rA \mid \rC, \rB, \rD) = \mi{\rA}{\rB \mid \rC, \rD}. \qed
	\end{align*}
	
\end{proof}

\begin{proposition}\label{prop:info-decrease}
  For random variables $\rA, \rB, \rC,\rD$, if $ \rA \perp \rD \mid \rB,\rC$, then, 
  \[\mi{\rA}{\rB \mid \rC} \geq \mi{\rA}{\rB \mid \rC, \rD}.\]
\end{proposition}
 \begin{proof}
 Since $\rA \perp \rD \mid \rB,\rC$, by \itfacts{cond-reduce}, $\HH(\rA \mid \rB,\rC) = \HH(\rA \mid \rB,\rC,\rD)$. Moreover, since conditioning can only reduce the entropy (again by \itfacts{cond-reduce}), 
  \begin{align*}
 	\mi{\rA}{\rB \mid  \rC} &= \HH(\rA \mid \rC) - \HH(\rA \mid \rB,\rC) \geq \HH(\rA \mid \rD,\rC) - \HH(\rA \mid \rB,\rC) \\
	&= \HH(\rA \mid \rD,\rC) - \HH(\rA \mid \rB,\rC,\rD) = \mi{\rA}{\rB \mid \rC,\rD}. \qed
 \end{align*}

\end{proof}


\subsection{Background on Communication and Information Complexity}\label{sec:cc-ic}

\paragraph{Communication complexity.} We briefly review the standard definitions of the two-party communication model of Yao~\cite{Yao79}.  See the text by Kushilevitz and Nisan~\cite{KushilevitzN97} for 
an extensive overview of communication complexity. In Section~\ref{sec:lower-bound-hpc}, we also use a standard generalization of this model to allow for more than two players, but we defer
the necessary definitions to that section. 

Let $P: \mathcal{X} \times \mathcal{Y} \rightarrow \mathcal{Z}$ be a relation.  Alice receives an input $X
\in \mathcal{X}$ and Bob receives $Y \in \mathcal{Y}$, where $(X,Y)$ are chosen from a
joint distribution $\cD$ over $\mathcal{X} \times \mathcal{Y}$. We allow players to have access to both public and private randomness. 
They communicate with each other by exchanging messages such that each message
depends only on the private input and random bits of the player sending the message, and the already communicated messages plus the public randomness. 
At the end, one of the players need to output an answer  $Z$ such that $Z \in P(X,Y)$.  

We use $\prot$ to denote a protocol used by the players. We always assume that the protocol $\prot$ can be randomized (using both public and
private randomness), \emph{even against a prior distribution $\cD$ of inputs}. For any
$0 < \delta < 1$, we say $\prot$ is a $\delta$-error protocol for $P$ over a distribution
$\cD$, if the probability that for an input $(X,Y)$, $\prot$ outputs some $Z$ where $Z \notin P(X,Y)$ is at most
$\delta$ (the probability is taken over the randomness of \emph{both} the distribution and the protocol).

\begin{definition}[Communication cost]
  The \emph{communication cost} of a protocol $\prot$ on an input
  distribution $\cD$, denoted by $\CC{\prot}{\cD}$, is the \emph{worst-case} bit-length of the transcript
  communicated between Alice and Bob in the protocol $\prot$, when the inputs are chosen from $\cD$.

\end{definition}

Communication complexity of a problem $P$ is defined as the minimum communication cost of a protocol $\prot$ that solves $P$ on every distribution $\cD$ with probability at least $2/3$. 



\paragraph{Information complexity.} There are several possible definitions of information
cost of a communication prtocol that have been considered depending on the application (see, e.g.,~\cite{ChakrabartiSWY01,Bar-YossefJKS02,BarakBCR10,BravermanR11,BravermanEOPV13}).  
We use the notion of \emph{internal information cost}~\cite{BarakBCR10} that measures the average amount of information each player learns about the input of the other player
by observing the transcript of the protocol. %Formally, 
\begin{definition}[Information cost]
  Consider an input distribution $\cD$ and a protocol $\prot$. Let $(\rX,\rY) \sim \cD$ denote the random variables for the input of Alice and Bob and $\rProt$ be the 
  the random variable for the transcript of the protocol \emph{concatenated} with the public randomness $\rR$ used by $\prot$. 
  The \emph{(internal) information cost} of $\prot$ with respect to
  $\cD$ is $\ICost{\prot}{\cD}:=\mii{\rProt}{\rX \mid \rY}{\cD} + \mii{\rProt}{\rY \mid \rX}{\cD}$. 

\end{definition}

One can also define information complexity of a problem $P$ similar to communication complexity with respect to the information cost. However, we avoid presenting this definition formally due to some subtle technical issues
 that need to be addressed which lead to multiple different but similar-in-spirit definitions. As such, we state our results directly in terms of information cost. 

%%
Note that any public coin protocol is a distribution over private coins protocols, run by
first using public randomness to sample a random string $\rR=R$ and then running the
corresponding private coin protocol $\prot^R$. We also use $\rProt^R$ to denote the transcript of the protocol $\prot^R$. 
We have the following standard proposition. 
\begin{proposition}\label{prop:public-random}
	 For any distribution $\cD$ and any protocol $\prot$ with public randomness $\bR$, $$\ICost{\prot}{\cD} = \mii{\rProt}{\rX \mid \rY, \rR}{\cD} + \mii{\rProt}{\rY \mid \rX, \rR}{\cD} = \E_{R \sim \rR}\bracket{\ICost{\prot^R}{\cD}}.$$ 
\end{proposition}
\begin{proof} By definition of internal information cost, 
\begin{align*}
	\ICost{\prot}{\cD} &= \mii{\rProt}{\rX \mid \rY}{\cD} + \mii{\rProt }{\rY \mid \rX}{\cD} = \mi{\rProt,\rR}{\rX \mid \rY} + \mi{\rProt,\rR}{\rY \mid \rX} \tag{$\Prot$ denotes the transcript and the public randomness} \\
	&= \mi{\rR}{\rX \mid \rY}  + \mi{\rProt}{\rX \mid \rY , \rR} + \mi{\rR}{\rY \mid \rX}  + \mi{\rProt}{\rY \mid \rX , \rR}  \tag{chain rule of mutual information, \itfacts{chain-rule}} \\
	&= \mi{\rProt}{\rX \mid \rY , \rR} + \mi{\rProt}{\rY \mid \rX , \rR} \tag{$\mi{\rR}{\rX \mid \rY}  = \mi{\rR}{\rY \mid \rX}  = 0$ since $\rR \perp \rX, \rY$ and \itfacts{info-zero}} \\
	&= \E_{R \sim \rR}\bracket{\mi{\rProt}{\rX \mid \rY , \rR=R} + \mi{\rProt}{\rY \mid \rX , \rR=R}} = \E_{R \sim \rR}\bracket{\ICost{\prot^R}{\cD}}, 
\end{align*}
concluding the proof.
\end{proof}


The following well-known proposition relates communication cost and information cost. 
\begin{proposition}[cf.~\cite{BravermanR11}]\label{prop:cc-ic}
  For any distribution $\cD$ and any protocol $\prot$: $\ICost{\prot}{\cD} \leq \CC{\prot}{\cD}$.
\end{proposition}
\begin{proof}
Let us assume first that $\prot$ only uses private randomness and thus $\rProt$ only contain the transcript.  
For any $b \in [\CC{\prot}{\cD}]$, we define $\Prot_b$ to be the $b$-th bit of the transcript. We have, 
\begin{align*}
	 \ICost{\prot}{\cD} &= \mi{\rProt}{\rX \mid \rY}+ \mi{\rProt }{\rY \mid \rX} \\
	 &= \sum_{b=1}^{\CC{\prot}{\cD}} \mi{\rProt_b}{\rX \mid \rProt^{<b},\rY} + \mi{\rProt_b}{\rY \mid \rProt^{<b},\rX} \tag{by chain rule of mutual information in \itfacts{chain-rule}}\\
	 &= \sum_{b=1}^{\CC{\prot}{\cD}} \E_{\Prot^{<b}}\bracket{\mi{\rProt_b}{\rX \mid \rProt^{<b} = \Prot^{<b},\rY} + \mi{\rProt_b}{\rY \mid \rProt^{<b}= \Prot^{<b},\rX}}.
\end{align*}
Consider each term in the RHS above. By conditioning on $\Prot^{<b}$, the player that transmit $\rProt_b$ would become fix. If this player is Alice, then $\mi{\rProt_b}{\rY \mid \rProt^{<b}= \Prot^{<b},\rX} = 0$, because
$\rProt_b$ is only a function of $(\rProt^{<b},\rX)$ in this case; similarly, if this player is Bob, then $\mi{\rProt_b}{\rX \mid \rProt^{<b}= \Prot^{<b},\rY} = 0$. Moreover, 
$\mi{\rProt_b}{\rX \mid \rProt^{<b} = \Prot^{<b},\rY} \leq \en{\rProt_b}  \leq 1$ and similarly $\mi{\rProt_b}{\rY \mid \rProt^{<b}= \Prot^{<b},\rX} \leq 1$. As such, the above term can be upper bounded 
by $\CC{\prot}{\cD}$. To finalize the proof, note that by Proposition~\ref{prop:public-random}, for any public-coin protocol $\prot$, $\ICost{\prot}{\cD} = \E_{R \sim \rR}\bracket{\ICost{\prot^R}{\cD}} \leq \E_{R \sim \rR}\bracket{\CC{\prot^R}{\cD}} \leq \CC{\prot}{\cD}$, where
the first inequality is by the first part of the argument.  
\end{proof}

Proposition~\ref{prop:cc-ic} provides a convinent way of proving communication complexity lower bounds by lower bounding information cost of any protocol. 


\subsection*{Rectangle Property of Communication Protocols}\label{sec:statstics}
We conclude this section by mentioning some basic properties of communication protocols. For any protocol $\prot$ and inputs $x \in \mathcal{X}$ and $y \in \mathcal{Y}$, we define $\Prot_{x,y}$ as the transcript
of the protocol conditioned on the input $x$ to Alice and input $y$ to Bob. Note that for randomized protocols, $\Prot_{x,y}$ is a random variable which we denote by $\rProt_{x,y}$. 

The following is referred to as the rectangle property of deterministic protocols. 
\begin{fact}[Rectangle property]\label{fact:rectangle}
	For any deterministic protocol $\prot$ and inputs $x,x' \in \mathcal{X}$ to Alice and $y,y' \in \mathcal{Y}$ to Bob, if $\Prot_{x,y} = \Prot_{x',y'}$, then $\Prot_{x,y'} = \Prot_{x',y}$. 
\end{fact}
\noindent
Fact~\ref{fact:rectangle} implies that the set of inputs consistent with any transcript $\Prot_{x,y}$ of a deterministic protocol forms a combinatorial rectangle. 
One can also extend the rectangle property of deterministic protocols to randomized protocols using the following fact. 
\begin{fact}[Cut-and-paste property; cf.~\cite{Bar-YossefJKS02}]\label{fact:r-rectangle}
	For any randomized protocol $\prot$ and inputs $x,x' \in \mathcal{X}$ to Alice and $y,y' \in \mathcal{Y}$ to Bob, $\hd{\rProt_{x,y}}{\rProt_{x',y'}} = \hd{\rProt_{x,y'}}{\rProt_{x',y}}$. 
\end{fact}

\newcommand{\etalchar}[1]{$^{#1}$}
\begin{thebibliography}{BvdBE{\etalchar{+}}22}

\bibitem[A17]{Assadi17sc}
Sepehr~Assadi A.
\newblock Tight space-approximation tradeoff for the multi-pass streaming set cover problem.
\newblock In {\em Proceedings of the 36th {ACM} {SIGMOD-SIGACT-SIGAI} Symposium on Principles of Database Systems, {PODS} 2017, Chicago, IL, USA, May 14-19, 2017}, pages 321--335, 2017.

\bibitem[A22]{Assadi22}
Sepehr~Assadi A.
\newblock A two-pass (conditional) lower bound for semi-streaming maximum matching.
\newblock In Joseph~(Seffi) Naor and Niv Buchbinder, editors, {\em Proceedings of the 2022 {ACM-SIAM} Symposium on Discrete Algorithms, {SODA} 2022, Virtual Conference / Alexandria, VA, USA, January 9 - 12, 2022}, pages 708--742. {SIAM}, 2022.

\bibitem[AA20]{AlonA20}
Noga Alon and Sepehr Assadi.
\newblock Palette sparsification beyond ({\(\Delta\)}+1) vertex coloring.
\newblock In Jaroslaw Byrka and Raghu Meka, editors, {\em Approximation, Randomization, and Combinatorial Optimization. Algorithms and Techniques, {APPROX/RANDOM} 2020, August 17-19, 2020, Virtual Conference}, volume 176 of {\em LIPIcs}, pages 6:1--6:22. Schloss Dagstuhl - Leibniz-Zentrum f{\"{u}}r Informatik, 2020.

\bibitem[ACK19]{AssadiCK19}
Sepehr Assadi, Yu~Chen, and Sanjeev Khanna.
\newblock Polynomial pass lower bounds for graph streaming algorithms.
\newblock In {\em {STOC}}, pages 265--276. {ACM}, 2019.

\bibitem[ACKP21]{AbboudCKP21}
Amir Abboud, Keren Censor{-}Hillel, Seri Khoury, and Ami Paz.
\newblock Smaller cuts, higher lower bounds.
\newblock {\em {ACM} Trans. Algorithms}, 17(4):30:1--30:40, 2021.

\bibitem[AD21]{AssadiD21}
Sepehr Assadi and Aditi Dudeja.
\newblock A simple semi-streaming algorithm for global minimum cuts.
\newblock In {\em {SOSA}}, pages 172--180. {SIAM}, 2021.

\bibitem[ADBV05]{Alvarez-HamelinDBV05}
J.~Ignacio Alvarez{-}Hamelin, Luca Dall'Asta, Alain Barrat, and Alessandro Vespignani.
\newblock Large scale networks fingerprinting and visualization using the k-core decomposition.
\newblock In {\em Advances in Neural Information Processing Systems 18 [Neural Information Processing Systems, {NIPS} 2005, December 5-8, 2005, Vancouver, British Columbia, Canada]}, pages 41--50, 2005.

\bibitem[AJJ{\etalchar{+}}22]{AssadiJJST22}
Sepehr Assadi, Arun Jambulapati, Yujia Jin, Aaron Sidford, and Kevin Tian.
\newblock Semi-streaming bipartite matching in fewer passes and optimal space.
\newblock In {\em {SODA}}, pages 627--669. {SIAM}, 2022.

\bibitem[AKSY20]{AssadiKSY20}
Sepehr Assadi, Gillat Kol, Raghuvansh Saxena, and Huacheng Yu.
\newblock Multi-pass graph streaming lower bounds for cycle counting, max-cut, matching size, and other problems.
\newblock In {\em 61st Annual {IEEE} Symposium on Foundations of Computer Science, {FOCS} (to appear)}, 2020.

\bibitem[AKZ22]{AssadiKZ22}
Sepehr Assadi, Gillat Kol, and Zhijun Zhang.
\newblock Rounds vs communication tradeoffs for maximal independent sets.
\newblock In {\em 63rd {IEEE} Annual Symposium on Foundations of Computer Science, {FOCS} 2022, Denver, CO, USA, October 31 - November 3, 2022}, pages 1193--1204. {IEEE}, 2022.

\bibitem[AN21]{AssadiN21}
Sepehr Assadi and Vishvajeet N.
\newblock Graph streaming lower bounds for parameter estimation and property testing via a streaming {XOR} lemma.
\newblock In Samir Khuller and Virginia~Vassilevska Williams, editors, {\em {STOC} '21: 53rd Annual {ACM} {SIGACT} Symposium on Theory of Computing, Virtual Event, Italy, June 21-25, 2021}, pages 612--625. {ACM}, 2021.

\bibitem[AR20]{AssadiR20}
Sepehr Assadi and Ran Raz.
\newblock Near-quadratic lower bounds for two-pass graph streaming algorithms.
\newblock In {\em {FOCS}}, pages 342--353. {IEEE}, 2020.

\bibitem[AYZ97]{AlonYZ97}
Noga Alon, Raphael Yuster, and Uri Zwick.
\newblock Finding and counting given length cycles.
\newblock {\em Algorithmica}, 17(3):209--223, 1997.

\bibitem[BBCR10]{BarakBCR10}
Boaz Barak, Mark Braverman, Xi~Chen, and Anup Rao.
\newblock How to compress interactive communication.
\newblock In {\em Proceedings of the 42nd {ACM} Symposium on Theory of Computing, {STOC} 2010, 5-8 June 2010}, pages 67--76, 2010.

\bibitem[BCD{\etalchar{+}}19]{BachrachCDELP19}
Nir Bachrach, Keren Censor{-}Hillel, Michal Dory, Yuval Efron, Dean Leitersdorf, and Ami Paz.
\newblock Hardness of distributed optimization.
\newblock In Peter Robinson and Faith Ellen, editors, {\em Proceedings of the 2019 {ACM} Symposium on Principles of Distributed Computing, {PODC} 2019, Toronto, ON, Canada, July 29 - August 2, 2019}, pages 238--247. {ACM}, 2019.

\bibitem[BCG20]{BeraCG19}
Suman~K. Bera, Amit Chakrabarti, and Prantar Ghosh.
\newblock Graph coloring via degeneracy in streaming and other space-conscious models.
\newblock In {\em 47th International Colloquium on Automata, Languages, and Programming, {ICALP} 2020, July 8-11, 2020, Saarbr{\"{u}}cken, Germany (Virtual Conference)}, pages 11:1--11:21, 2020.

\bibitem[BEO{\etalchar{+}}13]{BravermanEOPV13}
Mark Braverman, Faith Ellen, Rotem Oshman, Toniann Pitassi, and Vinod Vaikuntanathan.
\newblock A tight bound for set disjointness in the message-passing model.
\newblock In {\em 54th Annual {IEEE} Symposium on Foundations of Computer Science, {FOCS} 2013, 26-29 October, 2013, Berkeley, CA, {USA}}, pages 668--677, 2013.

\bibitem[BGKV14]{BonchiGKV14}
Francesco Bonchi, Francesco Gullo, Andreas Kaltenbrunner, and Yana Volkovich.
\newblock Core decomposition of uncertain graphs.
\newblock In Sofus~A. Macskassy, Claudia Perlich, Jure Leskovec, Wei Wang, and Rayid Ghani, editors, {\em The 20th {ACM} {SIGKDD} International Conference on Knowledge Discovery and Data Mining, {KDD} '14, New York, NY, {USA} - August 24 - 27, 2014}, pages 1316--1325. {ACM}, 2014.

\bibitem[BHNT15]{BhattacharyaHNT15}
Sayan Bhattacharya, Monika Henzinger, Danupon Nanongkai, and Charalampos~E. Tsourakakis.
\newblock Space- and time-efficient algorithm for maintaining dense subgraphs on one-pass dynamic streams.
\newblock In {\em Proceedings of the Forty-Seventh Annual {ACM} on Symposium on Theory of Computing, {STOC} 2015, Portland, OR, USA, June 14-17, 2015}, pages 173--182, 2015.

\bibitem[BJKS02]{Bar-YossefJKS02}
Ziv Bar{-}Yossef, T.~S. Jayram, Ravi Kumar, and D.~Sivakumar.
\newblock An information statistics approach to data stream and communication complexity.
\newblock In {\em 43rd Symposium on Foundations of Computer Science {(FOCS} 2002), 16-19 November 2002, Proceedings}, pages 209--218, 2002.

\bibitem[BKS02]{Bar-YossefKS02}
Ziv Bar{-}Yossef, Ravi Kumar, and D.~Sivakumar.
\newblock Reductions in streaming algorithms, with an application to counting triangles in graphs.
\newblock In {\em Proceedings of the Thirteenth Annual {ACM-SIAM} Symposium on Discrete Algorithms, January 6-8, 2002, San Francisco, CA, {USA.}}, pages 623--632, 2002.

\bibitem[BKV12]{BahmaniKV12}
Bahman Bahmani, Ravi Kumar, and Sergei Vassilvitskii.
\newblock Densest subgraph in streaming and mapreduce.
\newblock {\em {PVLDB}}, 5(5):454--465, 2012.

\bibitem[BR11]{BravermanR11}
Mark Braverman and Anup Rao.
\newblock Information equals amortized communication.
\newblock In {\em {IEEE} 52nd Annual Symposium on Foundations of Computer Science, {FOCS} 2011, October 22-25, 2011}, pages 748--757, 2011.

\bibitem[BvdBE{\etalchar{+}}22]{blikstadBEMN22}
Joakim Blikstad, Jan van~den Brand, Yuval Efron, Sagnik Mukhopadhyay, and Danupon Nanongkai.
\newblock Nearly optimal communication and query complexity of bipartite matching.
\newblock In {\em {FOCS}}, pages 1174--1185. {IEEE}, 2022.

\bibitem[CFHT20]{ChangFHT20}
Yi{-}Jun Chang, Martin Farach{-}Colton, Tsan{-}sheng Hsu, and Meng{-}Tsung Tsai.
\newblock Streaming complexity of spanning tree computation.
\newblock In {\em 37th International Symposium on Theoretical Aspects of Computer Science, {STACS} 2020, March 10-13, 2020, Montpellier, France}, pages 34:1--34:19, 2020.

\bibitem[CGMV20]{ChakrabartiGMV20}
Amit Chakrabarti, Prantar Ghosh, Andrew McGregor, and Sofya Vorotnikova.
\newblock Vertex ordering problems in directed graph streams.
\newblock In {\em Proceedings of the 2020 {ACM-SIAM} Symposium on Discrete Algorithms, {SODA} 2020, Salt Lake City, UT, USA, January 5-8, 2020}, pages 1786--1802, 2020.

\bibitem[Cha00]{Charikar00}
Moses Charikar.
\newblock Greedy approximation algorithms for finding dense components in a graph.
\newblock In Klaus Jansen and Samir Khuller, editors, {\em Approximation Algorithms for Combinatorial Optimization, Third International Workshop, {APPROX} 2000, Saarbr{\"{u}}cken, Germany, September 5-8, 2000, Proceedings}, volume 1913 of {\em Lecture Notes in Computer Science}, pages 84--95. Springer, 2000.

\bibitem[CKP{\etalchar{+}}21a]{ChenKPSSY21}
Lijie Chen, Gillat Kol, Dmitry Paramonov, Raghuvansh~R. Saxena, Zhao Song, and Huacheng Yu.
\newblock Almost optimal super-constant-pass streaming lower bounds for reachability.
\newblock In Samir Khuller and Virginia~Vassilevska Williams, editors, {\em {STOC} '21: 53rd Annual {ACM} {SIGACT} Symposium on Theory of Computing, Virtual Event, Italy, June 21-25, 2021}, pages 570--583. {ACM}, 2021.

\bibitem[CKP{\etalchar{+}}21b]{ChenKPSSY21b}
Lijie Chen, Gillat Kol, Dmitry Paramonov, Raghuvansh~R. Saxena, Zhao Song, and Huacheng Yu.
\newblock Near-optimal two-pass streaming algorithm for sampling random walks over directed graphs.
\newblock In Nikhil Bansal, Emanuela Merelli, and James Worrell, editors, {\em 48th International Colloquium on Automata, Languages, and Programming, {ICALP} 2021, July 12-16, 2021, Glasgow, Scotland (Virtual Conference)}, volume 198 of {\em LIPIcs}, pages 52:1--52:19. Schloss Dagstuhl - Leibniz-Zentrum f{\"{u}}r Informatik, 2021.

\bibitem[CKP{\etalchar{+}}23]{ChenKPSSY23}
Lijie Chen, Gillat Kol, Dmitry Paramonov, Raghuvansh~R. Saxena, Zhao Song, and Huacheng Yu.
\newblock Towards multi-pass streaming lower bounds for optimal approximation of max-cut.
\newblock In Nikhil Bansal and Viswanath Nagarajan, editors, {\em Proceedings of the 2023 {ACM-SIAM} Symposium on Discrete Algorithms, {SODA} 2023, Florence, Italy, January 22-25, 2023}, pages 878--924. {SIAM}, 2023.

\bibitem[CSWY01]{ChakrabartiSWY01}
Amit Chakrabarti, Yaoyun Shi, Anthony Wirth, and Andrew~Chi{-}Chih Yao.
\newblock Informational complexity and the direct sum problem for simultaneous message complexity.
\newblock In {\em 42nd Annual Symposium on Foundations of Computer Science, {FOCS} 2001, 14-17 October 2001}, pages 270--278, 2001.

\bibitem[CT06]{CoverT06}
Thomas~M. Cover and Joy~A. Thomas.
\newblock {\em Elements of information theory {(2.} ed.)}.
\newblock Wiley, 2006.

\bibitem[CZL{\etalchar{+}}20]{ChuZ00ZXZ20}
Deming Chu, Fan Zhang, Xuemin Lin, Wenjie Zhang, Ying Zhang, Yinglong Xia, and Chenyi Zhang.
\newblock Finding the best k in core decomposition: {A} time and space optimal solution.
\newblock In {\em 36th {IEEE} International Conference on Data Engineering, {ICDE} 2020, Dallas, TX, USA, April 20-24, 2020}, pages 685--696. {IEEE}, 2020.

\bibitem[DBS17]{DhulipalaBS17}
Laxman Dhulipala, Guy~E. Blelloch, and Julian Shun.
\newblock Julienne: {A} framework for parallel graph algorithms using work-efficient bucketing.
\newblock In Christian Scheideler and Mohammad~Taghi Hajiaghayi, editors, {\em Proceedings of the 29th {ACM} Symposium on Parallelism in Algorithms and Architectures, {SPAA} 2017, Washington DC, USA, July 24-26, 2017}, pages 293--304. {ACM}, 2017.

\bibitem[DBS18]{DhulipalaBS18}
Laxman Dhulipala, Guy~E. Blelloch, and Julian Shun.
\newblock Theoretically efficient parallel graph algorithms can be fast and scalable.
\newblock In Christian Scheideler and Jeremy~T. Fineman, editors, {\em Proceedings of the 30th on Symposium on Parallelism in Algorithms and Architectures, {SPAA} 2018, Vienna, Austria, July 16-18, 2018}, pages 393--404. {ACM}, 2018.

\bibitem[DNO19]{DobzinskiNO19}
Shahar Dobzinski, Noam Nisan, and Sigal Oren.
\newblock Economic efficiency requires interaction.
\newblock {\em Games Econ. Behav.}, 118:589--608, 2019.

\bibitem[EH66]{erdHos1966chromatic}
Paul Erd{\H{o}}s and Andr{\'a}s Hajnal.
\newblock On chromatic number of graphs and set-systems.
\newblock {\em Acta Math. Acad. Sci. Hungar}, 17(61-99):1, 1966.

\bibitem[ELM18]{EsfandiariLM18}
Hossein Esfandiari, Silvio Lattanzi, and Vahab~S. Mirrokni.
\newblock Parallel and streaming algorithms for k-core decomposition.
\newblock In Jennifer~G. Dy and Andreas Krause, editors, {\em Proceedings of the 35th International Conference on Machine Learning, {ICML} 2018, Stockholmsm{\"{a}}ssan, Stockholm, Sweden, July 10-15, 2018}, volume~80 of {\em Proceedings of Machine Learning Research}, pages 1396--1405. {PMLR}, 2018.

\bibitem[FHW12]{FrischknechtHW12}
Silvio Frischknecht, Stephan Holzer, and Roger Wattenhofer.
\newblock Networks cannot compute their diameter in sublinear time.
\newblock In Yuval Rabani, editor, {\em Proceedings of the Twenty-Third Annual {ACM-SIAM} Symposium on Discrete Algorithms, {SODA} 2012, Kyoto, Japan, January 17-19, 2012}, pages 1150--1162. {SIAM}, 2012.

\bibitem[FKM{\etalchar{+}}05]{FeigenbaumKMSZ05}
Joan Feigenbaum, Sampath Kannan, Andrew McGregor, Siddharth Suri, and Jian Zhang.
\newblock On graph problems in a semi-streaming model.
\newblock {\em Theor. Comput. Sci.}, 348(2-3):207--216, 2005.

\bibitem[FT14]{Farach-ColtonT14}
Martin Farach{-}Colton and Meng{-}Tsung Tsai.
\newblock Computing the degeneracy of large graphs.
\newblock In Alberto Pardo and Alfredo Viola, editors, {\em {LATIN} 2014: Theoretical Informatics - 11th Latin American Symposium, Montevideo, Uruguay, March 31 - April 4, 2014. Proceedings}, volume 8392 of {\em Lecture Notes in Computer Science}, pages 250--260. Springer, 2014.

\bibitem[FT16]{Farach-ColtonT16}
Martin Farach{-}Colton and Meng{-}Tsung Tsai.
\newblock Tight approximations of degeneracy in large graphs.
\newblock In {\em {LATIN} 2016: Theoretical Informatics - 12th Latin American Symposium, Ensenada, Mexico, April 11-15, 2016, Proceedings}, pages 429--440, 2016.

\bibitem[GBGL20]{GalimbertiBGL20}
Edoardo Galimberti, Francesco Bonchi, Francesco Gullo, and Tommaso Lanciano.
\newblock Core decomposition in multilayer networks: Theory, algorithms, and applications.
\newblock {\em {ACM} Trans. Knowl. Discov. Data}, 14(1):11:1--11:40, 2020.

\bibitem[GLM19]{GhaffariLM19}
Mohsen Ghaffari, Silvio Lattanzi, and Slobodan Mitrovic.
\newblock Improved parallel algorithms for density-based network clustering.
\newblock In Kamalika Chaudhuri and Ruslan Salakhutdinov, editors, {\em Proceedings of the 36th International Conference on Machine Learning, {ICML} 2019, 9-15 June 2019, Long Beach, California, {USA}}, volume~97 of {\em Proceedings of Machine Learning Research}, pages 2201--2210. {PMLR}, 2019.

\bibitem[GO13]{GuruswamiO13}
Venkatesan Guruswami and Krzysztof Onak.
\newblock Superlinear lower bounds for multipass graph processing.
\newblock In {\em Proceedings of the 28th Conference on Computational Complexity, {CCC} 2013, K.lo Alto, California, USA, 5-7 June, 2013}, pages 287--298, 2013.

\bibitem[HSSW12]{HalldorssonSSW12}
Magn{\'{u}}s~M. Halld{\'{o}}rsson, Xiaoming Sun, Mario Szegedy, and Chengu Wang.
\newblock Streaming and communication complexity of clique approximation.
\newblock In {\em Automata, Languages, and Programming - 39th International Colloquium, {ICALP} 2012, Warwick, UK, July 9-13, 2012, Proceedings, Part {I}}, pages 449--460, 2012.

\bibitem[IKL{\etalchar{+}}12]{IvanyosKLSW12}
G{\'{a}}bor Ivanyos, Hartmut Klauck, Troy Lee, Miklos Santha, and Ronald de~Wolf.
\newblock New bounds on the classical and quantum communication complexity of some graph properties.
\newblock In {\em {FSTTCS}}, volume~18 of {\em LIPIcs}, pages 148--159. Schloss Dagstuhl - Leibniz-Zentrum f{\"{u}}r Informatik, 2012.

\bibitem[JKS03]{JayramKS03}
T.~S. Jayram, Ravi Kumar, and D.~Sivakumar.
\newblock Two applications of information complexity.
\newblock In {\em {STOC}}, pages 673--682. {ACM}, 2003.

\bibitem[KBVT15]{KhaouidBST15}
Wissam Khaouid, Marina Barsky, S.~Venkatesh, and Alex Thomo.
\newblock K-core decomposition of large networks on a single {PC}.
\newblock {\em Proc. {VLDB} Endow.}, 9(1):13--23, 2015.

\bibitem[Kla00]{klauck2000quantum}
Hartmut Klauck.
\newblock On quantum and probabilistic communication: Las vegas and one-way protocols.
\newblock In {\em 32nd Annual ACM Symposium on Theory of Computing ({STOC})}, pages 644--651, 2000.

\bibitem[KN97]{KushilevitzN97}
Eyal Kushilevitz and Noam Nisan.
\newblock {\em Communication complexity}.
\newblock Cambridge University Press, 1997.

\bibitem[KPP16]{KopelowitzPP16}
Tsvi Kopelowitz, Seth Pettie, and Ely Porat.
\newblock Higher lower bounds from the 3sum conjecture.
\newblock In Robert Krauthgamer, editor, {\em Proceedings of the Twenty-Seventh Annual {ACM-SIAM} Symposium on Discrete Algorithms, {SODA} 2016, Arlington, VA, USA, January 10-12, 2016}, pages 1272--1287. {SIAM}, 2016.

\bibitem[KPSY23]{KolPSY23}
Gillat Kol, Dmitry Paramonov, Raghuvansh~R. Saxena, and Huacheng Yu.
\newblock Characterizing the multi-pass streaming complexity for solving boolean csps exactly.
\newblock In Yael~Tauman Kalai, editor, {\em 14th Innovations in Theoretical Computer Science Conference, {ITCS} 2023, January 10-13, 2023, MIT, Cambridge, Massachusetts, {USA}}, volume 251 of {\em LIPIcs}, pages 80:1--80:15. Schloss Dagstuhl - Leibniz-Zentrum f{\"{u}}r Informatik, 2023.

\bibitem[LJS19]{LiuJS19}
Yang~P. Liu, Arun Jambulapati, and Aaron Sidford.
\newblock Parallel reachability in almost linear work and square root depth.
\newblock In David Zuckerman, editor, {\em 60th {IEEE} Annual Symposium on Foundations of Computer Science, {FOCS} 2019, Baltimore, Maryland, USA, November 9-12, 2019}, pages 1664--1686. {IEEE} Computer Society, 2019.

\bibitem[LSY{\etalchar{+}}22]{LiuSYDS22}
Quanquan~C. Liu, Jessica Shi, Shangdi Yu, Laxman Dhulipala, and Julian Shun.
\newblock Parallel batch-dynamic algorithms for k-core decomposition and related graph problems.
\newblock In Kunal Agrawal and I{-}Ting~Angelina Lee, editors, {\em {SPAA} '22: 34th {ACM} Symposium on Parallelism in Algorithms and Architectures, Philadelphia, PA, USA, July 11 - 14, 2022}, pages 191--204. {ACM}, 2022.

\bibitem[LWSX18]{LiWSX18}
Chao Li, Li~Wang, Shiwen Sun, and Chengyi Xia.
\newblock Identification of influential spreaders based on classified neighbors in real-world complex networks.
\newblock {\em Appl. Math. Comput.}, 320:512--523, 2018.

\bibitem[LZZ{\etalchar{+}}19]{LiZZQZL19}
Conggai Li, Fan Zhang, Ying Zhang, Lu~Qin, Wenjie Zhang, and Xuemin Lin.
\newblock Efficient progressive minimum k-core search.
\newblock {\em Proc. {VLDB} Endow.}, 13(3):362--375, 2019.

\bibitem[MB83]{MatulaB83}
David~W. Matula and Leland~L. Beck.
\newblock Smallest-last ordering and clustering and graph coloring algorithms.
\newblock {\em J. {ACM}}, 30(3):417--427, 1983.

\bibitem[McG14]{McGregor14}
Andrew McGregor.
\newblock Graph stream algorithms: a survey.
\newblock {\em {SIGMOD} Rec.}, 43(1):9--20, 2014.

\bibitem[MN20]{MukhopadhyayN20}
Sagnik Mukhopadhyay and Danupon Nanongkai.
\newblock Weighted min-cut: sequential, cut-query, and streaming algorithms.
\newblock In {\em {STOC}}, pages 496--509. {ACM}, 2020.

\bibitem[MTVV15]{McGregorTVV15}
Andrew McGregor, David Tench, Sofya Vorotnikova, and Hoa~T. Vu.
\newblock Densest subgraph in dynamic graph streams.
\newblock In {\em Mathematical Foundations of Computer Science 2015 - 40th International Symposium, {MFCS} 2015, Milan, Italy, August 24-28, 2015, Proceedings, Part {II}}, pages 472--482, 2015.

\bibitem[NW91]{NisanW91}
Noam Nisan and Avi Wigderson.
\newblock Rounds in communication complexity revisited.
\newblock In {\em Proceedings of the 23rd Annual {ACM} Symposium on Theory of Computing, May 5-8, 1991, New Orleans, Louisiana, {USA}}, pages 419--429, 1991.

\bibitem[NW93]{NisanW93}
Noam Nisan and Avi Wigderson.
\newblock Rounds in communication complexity revisited.
\newblock {\em {SIAM} J. Comput.}, 22(1):211--219, 1993.

\bibitem[PRV99]{PonzioRV99}
Stephen Ponzio, Jaikumar Radhakrishnan, and Srinivasan Venkatesh.
\newblock The communication complexity of pointer chasing: Applications of entropy and sampling.
\newblock In {\em Proceedings of the Thirty-First Annual {ACM} Symposium on Theory of Computing, May 1-4, 1999, Atlanta, Georgia, {USA}}, pages 602--611, 1999.

\bibitem[PS84]{PapadimitriouS84}
Christos~H. Papadimitriou and Michael Sipser.
\newblock Communication complexity.
\newblock {\em J. Comput. Syst. Sci.}, 28(2):260--269, 1984.

\bibitem[Raz92]{Razborov92}
Alexander~A. Razborov.
\newblock On the distributional complexity of disjointness.
\newblock {\em Theor. Comput. Sci.}, 106(2):385--390, 1992.

\bibitem[RSW18]{RubinsteinSW18}
Aviad Rubinstein, Tselil Schramm, and S.~Matthew Weinberg.
\newblock Computing exact minimum cuts without knowing the graph.
\newblock In {\em {ITCS}}, volume~94 of {\em LIPIcs}, pages 39:1--39:16. Schloss Dagstuhl - Leibniz-Zentrum f{\"{u}}r Informatik, 2018.

\bibitem[RW13]{RodittyW13}
Liam Roditty and Virginia~Vassilevska Williams.
\newblock Fast approximation algorithms for the diameter and radius of sparse graphs.
\newblock In Dan Boneh, Tim Roughgarden, and Joan Feigenbaum, editors, {\em Symposium on Theory of Computing Conference, STOC'13, Palo Alto, CA, USA, June 1-4, 2013}, pages 515--524. {ACM}, 2013.

\bibitem[SGJ{\etalchar{+}}13]{SariyuceGJWC13}
Ahmet~Erdem Sariy{\"{u}}ce, Bugra Gedik, Gabriela Jacques{-}Silva, Kun{-}Lung Wu, and {\"{U}}mit~V. {\c{C}}ataly{\"{u}}rek.
\newblock Streaming algorithms for k-core decomposition.
\newblock {\em Proc. {VLDB} Endow.}, 6(6):433--444, 2013.

\bibitem[Yao79]{Yao79}
Andrew~Chi{-}Chih Yao.
\newblock Some complexity questions related to distributive computing (preliminary report).
\newblock In {\em Proceedings of the 11h Annual {ACM} Symposium on Theory of Computing, April 30 - May 2, 1979, Atlanta, Georgia, {USA}}, pages 209--213, 1979.

\bibitem[Yeh20]{Yehudayoff20}
Amir Yehudayoff.
\newblock Pointer chasing via triangular discrimination.
\newblock {\em Comb. Probab. Comput.}, 29(4):485--494, 2020.

\end{thebibliography}
\end{document}
